% !TeX root = ../../Book3.tex
% 纲要标题：## 第13章　全局椭圆理论与非线性分析
% 纲要位置：几何学计划纲要.md，第 2608 行。
% 正文状态：初稿；尚未正式审核。

% 纲要内容开始
% 
% > 本章目标与证明范围：从局部估计和紧嵌入推出闭值域、有限核余核，再构造 Green 算子及椭圆复形的调和模型。第13.5—13.6节独立建立非线性共同工具，供 Yamabe、Einstein、几何流、伪全纯方程与形变理论调用；只学线性椭圆与 Hodge 的读者可以先跳过这两节。
% 
% 纲要内容结束

\chapter{全局椭圆理论与非线性分析}
\label{g3:ch:13}
\BookChapterNumber{g3:ch:13}

\begin{chapterabstract}
局部正则性说明椭圆方程的解有多光滑，全局理论则判断哪些右端可解、解有多少自由度以及怎样连续选择解。本章在闭流形上由局部估计与 Rellich 紧性证明核有限维、值域闭，并用形式伴随识别余核；随后构造 Green 算子，证明非负自伴椭圆算子的离散谱及相应的 Sobolev 展开。椭圆复形把这些分析结果组织成正合、余正合与调和分解。后两节从 Banach 微分、定量隐函数方法与有限维约化出发，建立 Sobolev 映射空间和非线性微分算子的实际运算规则，为后续几何方程提供可直接核对的求解条件。
\end{chapterabstract}

\begin{learninggoals}
\begin{itemize}
	\item 用有限覆盖拼接内部估计，并由紧嵌入证明核有限维及核正交补上的加强估计。
	\item 区分形式伴随与带定义域的 Hilbert 伴随，写出完整的 Fredholm 可解条件。
	\item 构造 Green 广义逆，证明紧正自伴谱分解并用特征系数刻画正则性。
	\item 检验微分复形的主符号正合，推导调和分解与有限维上同调代表。
	\item 证明定量逆函数、隐函数和分裂满射定理，计算有限维障碍映射。
	\item 为 Sobolev 映射选取可用的指数和坐标，核对乘法、复合及微分损失。
	\item 在明确的系数与初始正则性下实施非线性自举，辨别约束极小与对称性造成的自由度。
\end{itemize}
\end{learninggoals}

在圆周上求解 \(u'=f\)，局部积分总能给出原函数，但周期性要求 \(\int_{S^1}f=0\)。这个额外条件只占一个维度；满足它以后，解仍可加任意常数。两种有限维自由度分别属于右端和未知量，全局方程必须同时处理它们。椭圆正则性负责控制局部导数，紧流形上的整体拓扑与算子结构则决定哪些方向留下核和可解障碍。

\ProseParagraphBreak

把方程改成实值非线性方程 \(u'+u^2=f\)，积分条件变为
\[
	\int_{S^1}u^2=\int_{S^1}f.
\]
例如 \(f=0\) 时只有零解；但在零解处线性化仍是 \(\partial_x\)，其核含全部常数。线性化的核说明一阶方程未限制哪些方向，却不能独自保证这些方向延伸成真实解族。要判断解的局部形状，需要先反演可逆方向，再检验核方向上剩余的非线性方程。这正是后面有限维约化所保留的信息。

\ProseParagraphBreak

本章前四节沿着全局估计、闭值域、Green 算子和椭圆复形展开。闭流形的有限覆盖将第十二章的内部估计拼起来，第十一章的紧嵌入排除无限多个相互分离的归一化齐次解。伴随方程把余核转化为光滑截面，因而可解性由有限个积分条件决定。对非负自伴算子，逆的紧性进一步给出离散特征展开；对微分复形，主符号正合则把这个分析构造变成上同调的调和模型。

\ProseParagraphBreak

第13.5—13.6节建立非线性共同工具。逆函数、分裂满射与 Fredholm 约化使用不同的线性化条件，Sobolev 复合则同时取决于源维数、导数阶数与可积指数。只学习线性椭圆与 Hodge 理论的读者可先读前四节，待使用非线性几何方程时再回到后两节。所需函数空间、局部估计和紧性均来自前两章；紧算子的基础与谱分解可对照 Brezis 第6.2、6.4节\cite[第6.2、6.4节]{Brezis2011}。边界正则性与边值条件放在下一章，闭流形的分析结论在本章完整建立。

% !TeX root = ../../Book3.tex
% 纲要标题：### 13.1 紧流形上的全局估计
% 纲要位置：几何学计划纲要.md，第 2612 行。
% 正文状态：初稿；尚未正式审核。

% 纲要内容开始
% 
% * 有限覆盖与局部估计的粘合
% * 含低阶范数的先验不等式
% * 核的有限维性
% * 核的正交补上的先验估计
% 
% 纲要内容结束

\section{紧流形上的全局估计}
\label{g3:sec:1301}

局部估计控制一个坐标片内的导数。要解整个流形上的方程，还须知道这些估计能否统一、齐次方程有多少自由度，以及去掉这些自由度后能否稳定地反解。本节在闭流形上把三个问题接起来：有限覆盖给出全局估计，紧嵌入限制核的维数，正交归一化去掉估计中最后的零阶项。

\ProseParagraphBreak

这里的闭流形指紧且无边界的光滑流形。固定光滑底度量、正体积密度、有限秩丛的度量与联络；复 Hermitian 内积仍对第一变量线性，\((u,v)_2=\int_M\langle u,v\rangle\mu_g\)。算子的阶数记为 \(d\)，流形维数记为 \(n\)。

\subsection{有限覆盖与局部估计的粘合}
\label{g3:subsec:1301:01}
先说明 Sobolev 空间与微分算子的具体对应。在有限平凡化中，光滑 \(d\) 阶算子有表达式
\[
	P=\sum_{|\alpha|\le d}A_\alpha(x)\mathrm D^\alpha,
	\qquad \mathrm D_j=-\mathrm i\partial_j.
\]
对每个整数 \(k\ge0\)，乘积求导至 \(k\) 阶时，\(Pu\) 只涉及 \(u\) 的至 \(k+d\) 阶导数及有限阶系数导数。紧集上的系数界与\cref{g3:prop:sobolev-norm-equivalence}因此给出有界映射
\[
	P:W^{k+d,p}(M;E)\longrightarrow W^{k,p}(M;F),\qquad1\le p<\infty.
\]
这个作用与分布意义的作用一致，因为光滑逼近可代入每个测试配对。椭圆性将提供反方向的估计。

\begin{theorem}[全局椭圆估计]\label{g3:thm:global-elliptic-hk}
	设 \(M^n\) 为非空闭光滑流形，\(n\ge1\)，\(E,F\to M\) 为同秩有限秩光滑实或复向量丛，固定光滑度量与联络。设 \(P:C^\infty(E)\to C^\infty(F)\) 为光滑系数的椭圆微分算子，阶数 \(d\ge1\)。对每个整数 \(k\ge0\)，存在 \(C_k\)，使
	\begin{equation}\label{g3:eq:global-elliptic-hk}
		\|u\|_{H^{k+d}}\le C_k\bigl(\|Pu\|_{H^k}+\|u\|_2\bigr)
	\end{equation}
	对所有 \(u\in H^{k+d}(E)\) 成立。若分布截面 \(u\) 满足 \(Pu\in H^k(F)\)，则 \(u\in H^{k+d}(E)\)，并满足同一估计。
\end{theorem}
\begin{proof}
	选有限平凡化 \(U_a\)，从属单位分解 \(\chi_a\)，以及 \(\psi_a\in C_c^\infty(U_a)\)，使 \(\psi_a=1\) 于 \(\operatorname{supp}\chi_a\) 附近。第十二章的分布估计\cref{g3:thm:elliptic-distribution-l2}在每个坐标片中取 \(s=k,t=0\)，得到
	\[
		\|(\chi_a u)_a\|_{H^{k+d}(\mathbb R^n)}
		\le C_a\bigl(\|(\psi_a Pu)_a\|_{H^k(\mathbb R^n)}+\|(\psi_a u)_a\|_2\bigr).
	\]
	下标表示相应丛的局部分量。局部定理已经控制了截断交换子；这一步无需把 \(P(\chi_a u)\) 错写成 \(\chi_a Pu\)。对有限个 \(a\) 求和，使用 Sobolev 范数等价及光滑乘法的有界性，即得式\eqref{g3:eq:global-elliptic-hk}。光滑稠密性使估计延伸到 \(H^{k+d}\)。

	对任意分布 \(u\)，先在每张坐标片使用同一定理的正则性部分，得到 \(u\in H^{k+d}_{\mathrm{loc}}\)。紧性与有限覆盖使其成为全局 \(H^{k+d}\) 截面，再应用已证估计。这个次序并未预先要求未知分布具有 \(L^2\) 范数。
\end{proof}

闭流形上的每个点都位于一个内部坐标片中。若 \(M\) 有边界，覆盖边界的半球不是第十二章的内部区域；边界条件和相应估计将在下一章另行建立。对变化的几何族，式\eqref{g3:eq:global-elliptic-hk}中的常数仍依赖覆盖、截断、度量、标架和所需系数界，固定流形上的有限性不自动给出族上的一致性。

\subsection{含低阶范数的先验不等式}
\label{g3:subsec:1301:02}
零阶项记录了尚未由方程控制的低频部分。要在其他可积指数中得到相同形式，先把低阶导数压到任意小倍数的高阶范数中。

\begin{lemma}\label{g3:lem:compact-sobolev-interpolation}
	设 \(M\) 为闭光滑流形，\(E\) 为光滑有限秩度量丛，\(q\ge1\) 为整数，\(1<p<\infty\)。对每个 \(\varepsilon>0\)，存在 \(C_\varepsilon\)，使
	\[
		\|u\|_{W^{q-1,p}}\le\varepsilon\|u\|_{W^{q,p}}+C_\varepsilon\|u\|_{L^p}
		\qquad(u\in W^{q,p}(E)).
	\]
\end{lemma}
\begin{proof}
	\(q=1\) 时取 \(C_\varepsilon=1\) 即可。若某个 \(q>1\)、\(\varepsilon>0\) 的结论不成立，归一化后可取 \(u_j\)，使
	\[
		\|u_j\|_{W^{q-1,p}}=1,\qquad
		\varepsilon\|u_j\|_{W^{q,p}}+j\|u_j\|_p<1.
	\]
	由\cref{g3:thm:compact-bundle-rellich}，一个子列在 \(W^{q-1,p}\) 中收敛。它的 \(L^p\) 范数趋零，故极限只能为零，却又应有 \(W^{q-1,p}\) 范数一，矛盾。
\end{proof}

\begin{theorem}\label{g3:thm:global-elliptic-wkp}
	设 \(M^n\) 为非空闭光滑流形，\(n\ge1\)，\(E,F\) 为同秩有限秩光滑度量丛，\(P:C^\infty(E)\to C^\infty(F)\) 是光滑系数、阶 \(d\ge1\) 的椭圆算子。固定整数 \(k\ge0\) 及 \(1<p<\infty\)。则
	\begin{equation}\label{g3:eq:global-elliptic-wkp}
		\|u\|_{W^{k+d,p}}\le C\bigl(\|Pu\|_{W^{k,p}}+\|u\|_p\bigr)
	\end{equation}
	对 \(u\in W^{k+d,p}(E)\) 成立；若分布截面满足 \(Pu\in W^{k,p}(F)\)，则它属于 \(W^{k+d,p}(E)\)，并满足此式。
\end{theorem}
\begin{proof}
	\(k=0\) 时，在一组内外嵌套坐标片上应用\cref{g3:thm:local-elliptic-lp-apriori}，有限相加即得结论。

	为处理高阶，先令 \(u\) 光滑。在每个稍大的坐标片中，对全部 \(|\beta|=k\) 写出
	\[
		P(\partial^\beta u)=\partial^\beta(Pu)+[P,\partial^\beta]u.
	\]
	交换子中至少一次导数落在系数上，所以其作用于 \(u\) 的阶数至多为 \(k+d-1\)。局部 \(W^{d,p}\) 估计用于 \(\partial^\beta u\)，再加上较低阶导数并以有限片拼接，得到
	\[
		\|u\|_{W^{k+d,p}}
		\le C\bigl(\|Pu\|_{W^{k,p}}+\|u\|_{W^{k+d-1,p}}\bigr).
	\]
	坐标、联络与截断带来的项也都包含在后一个范数中。上一引理取 \(q=k+d\)，选择 \(\varepsilon\) 使 \(C\varepsilon<1/2\)，移项得到所需估计。光滑稠密性处理一般 \(W^{k+d,p}\) 截面。

	分布解先由\cref{g3:thm:elliptic-distribution-lp}逐片提升为 \(W^{k+d,p}\)，再用紧性拼合并代入估计。实丛可复化后取实部，结论相同。
\end{proof}

\begin{example}\label{g3:exm:global-zero-order-term}
	连通闭流形上的非负 Laplace 算子 \(\Delta_g\) 满足
	\(\|u\|_{H^2}\le C(\|\Delta_gu\|_2+\|u\|_2)\)。常数函数的 Laplace 为零，故尚未规定均值时不能删除零阶项。相反，\(P=1+\Delta_g\) 的齐次解与自身配对给出 \(\|u\|_2^2+\|\mathrm du\|_2^2=0\)，其核为零；下面的估计将说明此时可以去掉该项。
\end{example}

\subsection{核的有限维性}
\label{g3:subsec:1301:03}
令 \(\mathcal N_P\) 表示齐次方程的光滑解空间。局部正则性使这个空间也包含全部分布齐次解，因此核的大小由全局几何决定。

\begin{proposition}\label{g3:prop:finite-elliptic-kernel}
	设 \(M^n\) 为非空闭光滑流形，\(n\ge1\)，\(E,F\) 为同秩光滑有限秩度量丛，\(P:C^\infty(E)\to C^\infty(F)\) 为光滑椭圆算子，阶 \(d\ge1\)。则
	\[
		\mathcal N_P=\{u\in\mathcal D'(E):Pu=0\}
		\subset C^\infty(E)
	\]
	为有限维空间。它是 \(L^2(E)\) 的闭子空间，所有 Sobolev 实现的核都等于它。
\end{proposition}
\begin{proof}
	分布解光滑性来自\cref{g3:cor:elliptic-smoothness}。若 \(u_j\in\mathcal N_P\) 在 \(L^2\) 中趋于 \(u\)，则对任意光滑测试截面，分部积分并取极限可得 \(Pu=0\) 分布成立；所以 \(u\in\mathcal N_P\)，核闭。

	若核无限维，逐次取与前面有限多个向量正交的单位向量，可构造 \(L^2\) 正交归一列 \(e_j\in\mathcal N_P\)。全局估计给出 \(\|e_j\|_{H^d}\le C\)。由 \(d\ge1\) 及 Rellich 紧性，一个子列在 \(L^2\) 中收敛；但不同项的 \(L^2\) 距离恒为 \(\sqrt2\)，矛盾。因此核有限维。任何 Sobolev 核都包含于分布核，而每个光滑核元素属于全部这些空间，故各实现的核相同。
\end{proof}

取 \(\mathcal N_P\) 的 \(L^2\) 正交归一基 \(e_1,\ldots,e_r\)，正交投影为
\begin{equation}\label{g3:eq:elliptic-kernel-projection}
	\Pi_Pu=\sum_{j=1}^r(u,e_j)_2e_j.
\end{equation}
\symidx{\(\mathcal N_P,\Pi_P\)：椭圆算子的光滑核及其 \(L^2\) 正交投影}
它与基的选择无关，因为 \(u-\Pi_Pu\) 与核正交。每个 \(e_j\) 光滑，Hölder 不等式还给出
\[
	\|\Pi_Pu\|_{W^{q,p}}\le C_{q,p}\|u\|_{L^p},
	\qquad q\ge0,\quad1<p<\infty.
\]
因此式\eqref{g3:eq:elliptic-kernel-projection}也定义了 \(L^p\) 上的有界有限秩投影；即使 \(p\ne2\)，我们仍用这些积分配对规定“与核正交”。

\subsection{核的正交补上的先验估计}
\label{g3:subsec:1301:04}
齐次解可以任意加到一个特解上。规定特解与核正交，既选出唯一候选，也恢复稳定估计。

\begin{theorem}\label{g3:thm:kernel-orthogonal-estimate}
	设 \(M^n\) 为非空闭光滑流形，\(n\ge1\)，\(E,F\) 为同秩光滑有限秩度量丛，\(P:C^\infty(E)\to C^\infty(F)\) 为光滑椭圆 \(d\) 阶算子，\(d\ge1\)。对整数 \(k\ge0\)、\(1<p<\infty\)，存在 \(C_{k,p}\)，使
	\begin{equation}\label{g3:eq:kernel-orthogonal-estimate}
		\|u\|_{W^{k+d,p}}\le C_{k,p}\|Pu\|_{W^{k,p}}
	\end{equation}
	对所有 \(u\in W^{k+d,p}(E)\)、\(\Pi_Pu=0\) 成立。特别地，任意 \(u\in W^{k+d,p}(E)\) 满足同一右端控制 \(\|u-\Pi_Pu\|_{W^{k+d,p}}\) 的估计。
\end{theorem}
\begin{proof}
	先证明 \(\Pi_Pu=0\) 时 \(\|u\|_p\le C\|Pu\|_p\)。若不成立，可取 \(u_j\in W^{d,p}\)，使
	\[
		\Pi_Pu_j=0,\qquad \|u_j\|_p=1,\qquad \|Pu_j\|_p\longrightarrow0.
	\]
	式\eqref{g3:eq:global-elliptic-wkp}使它们在 \(W^{d,p}\) 中有界。由 Rellich 紧性，可取子列在 \(L^p\) 中趋于 \(u\)。于是 \(\|u\|_p=1\)、\(\Pi_Pu=0\)，并且测试配对中 \(Pu=0\)。分布正则性使 \(u\in\mathcal N_P\)，所以 \(u=\Pi_Pu=0\)，矛盾。

	现在把 \(\|u\|_p\le C\|Pu\|_p\) 代入全局 \(W^{k+d,p}\) 估计，即得式\eqref{g3:eq:kernel-orthogonal-estimate}。最后 \(P\Pi_Pu=0\)，所以将同一估计用于 \(u-\Pi_Pu\) 即可。
\end{proof}

这个矛盾论证也说明常数的含义：若归一化截面能让方程残差越来越小，紧性会产生一个非零齐次解。正交条件恰好排除了这个极限。因而当一族算子的核发生跳跃时，补空间上的反解常数可能变大；指数的稳定性与这种常数的一致性是两个不同的问题。
% 纲要条目排为编号小节标题，并进入目录与书签；说明段仍保留于纲要注释中。

% !TeX root = ../../Book3.tex
% 纲要标题：### 13.2 闭值域与 Fredholm 择一定理
% 纲要位置：几何学计划纲要.md，第 2619 行。
% 正文状态：初稿；尚未正式审核。

% 纲要内容开始
% 
% * 值域的闭性
% * 形式伴随的正则性
% * 余核与伴随核
% * 可解条件和 Fredholm 指数
% * 指数对连续扰动的稳定性
% 
% 纲要内容结束

\section{闭值域与 Fredholm 择一定理}
\label{g3:sec:1302}

有了核上的投影与补空间估计，方程能否求解就转化为右端是否落在值域中。闭值域保证近似可解的右端确实可解；形式伴随把所有障碍写成有限个积分条件。这是本节的 Fredholm 选择。\footnote{Erik Ivar Fredholm（1866--1927），瑞典数学家，研究积分方程、谱理论与数学物理。}它把矩阵方程中“右端与转置核正交”的条件推广到闭流形上的椭圆方程。

\subsection{值域的闭性}
\label{g3:subsec:1302:01}
若 \(Pu_j\) 收敛，\(u_j\) 未必收敛，因为可以给它任意加上核元素。先减去核投影，上一节的估计就把右端的 Cauchy 性传给解。

\begin{proposition}\label{g3:prop:elliptic-closed-range}
	设 \(M^n\) 为非空闭光滑流形，\(n\ge1\)，\(E,F\) 为同秩光滑有限秩度量丛，\(P:C^\infty(E)\to C^\infty(F)\) 是光滑椭圆算子，阶 \(d\ge1\)。固定底流形的光滑黎曼度量以及各丛的光滑参考联络，Sobolev范数用这些数据定义；积分与形式伴随采用该度量的体积密度及给定的纤维度量。记 \(\Pi_P\) 为到光滑核 \(\ker P\) 的 \(L^2\) 正交投影。则 \(P:H^d(E)\to L^2(F)\) 的值域 \(\mathcal R_P\) 闭，且
	\[
		P:\{u\in H^d(E):\Pi_Pu=0\}\longrightarrow\mathcal R_P
	\]
	为双射，其逆有界。
\end{proposition}
\begin{proof}
	设 \(f_j=Pu_j\to f\) 于 \(L^2\)，令 \(v_j=u_j-\Pi_Pu_j\)。\Cref{g3:thm:kernel-orthogonal-estimate}给出
	\[
		\|v_j-v_\ell\|_{H^d}\le C\|f_j-f_\ell\|_2.
	\]
	故 \(v_j\) 在完备空间 \(H^d\) 中收敛到 \(v\)。算子的有界性给出 \(Pv=f\)，投影的连续性给出 \(\Pi_Pv=0\)，所以 \(f\in\mathcal R_P\)。核正交截面中的齐次解为零，故所述映射单射；减去核投影又不改变右端，故它满射。逆的范数界正是上一节的估计。
\end{proof}

证明中的关键下界来自上一节的矛盾紧性论证。闭值域使残差趋零具有求解意义：只要右端 \(f\) 可以被某列 \(Pu_j\) 在 \(L^2\) 中逼近，就能在去掉核以后得到一个真正满足 \(Pu=f\) 的极限。

\begin{example}\label{g3:exm:dense-nonclosed-range}
	在 \(\ell^2\) 中令 \(T(x_j)=(x_j/j)\)。它有界且单射，值域包含所有有限支序列，因而稠密。令 \(y=(1/j)_{j\ge1}\in\ell^2\)，其有限截断均在值域中，且收敛到 \(y\)；但 \(Tx=y\) 要求 \(x_j=1\)，不属于 \(\ell^2\)，故值域不闭。这里伴随核为零，却仍有不可解右端，说明正交条件要成为充分条件，闭值域不能缺少。
\end{example}

\subsection{形式伴随的正则性}
\label{g3:subsec:1302:02}
第十二章的形式伴随 \(P^*\) 由积分恒等式确定：
\[
	(Pu,v)_2=(u,P^*v)_2,\qquad u\in C^\infty(E),\quad v\in C^\infty(F).
\]
闭流形没有边界项。按\cref{g3:prop:symbol-calculus}，其主符号为 \(\sigma_d(P^*)(x,\xi)=\sigma_d(P)(x,\xi)^*\)，故 \(P^*\) 同样椭圆，阶数仍为 \(d\)。于是 \(P^*v=0\) 的分布解也是光滑的，且 \(\mathcal N_{P^*}\) 有限维。

\ProseParagraphBreak

形式伴随是微分表达式。若把 \(P\) 看成从 \(H^d\) 到 \(L^2\) 的有界映射，它关于这两个 Hilbert 范数的有界伴随取值于 \(H^d\)，并不是同一个表达式。谱论需要的是把 \(P\) 看成 \(L^2\) 上定义域为 \(H^d\) 的算子，下面把这一点准确落实。

\begin{proposition}\label{g3:prop:elliptic-adjoint-domain}
	设 \(M^n\) 为非空闭光滑流形，\(n\ge1\)，\(E,F\) 为同秩有限秩光滑 Hermitian 丛，\(P:C^\infty(E)\to C^\infty(F)\) 为光滑椭圆 \(d\) 阶算子，\(d\ge1\)。固定底流形的光滑黎曼度量以及各丛的光滑参考联络，Sobolev范数用这些数据定义；积分与形式伴随采用该度量的体积密度及给定的纤维度量。以 \(P^*\) 表示由固定体积密度及丛度量确定的形式伴随。把 \(\mathsf P\) 定义为 \(L^2(E)\to L^2(F)\) 的算子，域为 \(H^d(E)\)，作用为 \(P\)。则 \(\mathsf P\) 稠定且闭。其 Hilbert 伴随的域恰为 \(H^d(F)\)，在此域上作用为 \(P^*\)。
\end{proposition}
\begin{proof}
	光滑截面在 \(L^2\) 中稠密，故定义域稠密。若 \(u_j\to u\) 于 \(L^2(E)\)、\(Pu_j\to f\) 于 \(L^2(F)\)，全局估计给出
	\[
		\|u_j-u_\ell\|_{H^d}
		\le C\bigl(\|Pu_j-Pu_\ell\|_2+\|u_j-u_\ell\|_2\bigr).
	\]
	因此 \(u_j\) 在 \(H^d\) 中收敛，极限在 \(L^2\) 中只能是 \(u\)，且 \(Pu=f\)。这正是图闭的含义。

	按定义，\(v\in L^2(F)\) 属于 Hilbert 伴随的域，等价于存在 \(w\in L^2(E)\)，使
	\((Pu,v)_2=(u,w)_2\) 对全部 \(u\in H^d(E)\) 成立。对光滑 \(u\) 作测试，得 \(P^*v=w\) 分布成立。伴随的椭圆正则性使 \(v\in H^d(F)\)。反之，若 \(v\in H^d(F)\)，先对光滑 \(u,v\) 分部积分，再由 \(H^d\) 光滑逼近取极限，得到上述恒等式，其中 \(w=P^*v\in L^2\)。故两种定义域恰好相同。
\end{proof}

\subsection{余核与伴随核}
\label{g3:subsec:1302:03}
对右端空间中的闭子空间 \(\mathcal R_P\)，正交补中的向量记录了全部线性阻碍。形式伴随的正则性把这些原本只有 \(L^2\) 意义的向量变成光滑截面。

\termidx[Fredholm xuanze]{Fredholm 选择}
\foreignidx{Fredholm alternative}
\begin{theorem}[椭圆 Fredholm 选择]\label{g3:thm:elliptic-fredholm}
	设 \(M^n\) 为非空闭光滑流形，\(n\ge1\)，\(E,F\) 为同秩光滑有限秩度量丛，\(P:C^\infty(E)\to C^\infty(F)\) 为光滑椭圆算子，阶 \(d\ge1\)。固定底流形的光滑黎曼度量以及各丛的光滑参考联络，Sobolev范数用这些数据定义；积分与形式伴随采用该度量的体积密度及给定的纤维度量。记 \(P^*\) 为形式伴随，\(\mathcal N_Q\) 为 \(Q=P,P^*\) 在光滑截面上的核，\(\Pi_P\) 为到 \(\mathcal N_P\) 的 \(L^2\) 正交投影。则
	\begin{equation}\label{g3:eq:elliptic-range-orthogonal}
		\operatorname{ran}(P:H^d(E)\to L^2(F))
		=\mathcal N_{P^*}^{\perp}.
	\end{equation}
	因此 \(L^2(F)/\operatorname{ran}P\) 自然同构于 \(\mathcal N_{P^*}\)，特别地余维有限。对 \(f\in L^2(F)\)，方程 \(Pu=f\) 可解当且仅当
	\[
		(f,h)_2=0\qquad\text{对全部 }h\in\mathcal N_{P^*}.
	\]
	可解时有唯一满足 \(\Pi_Pu=0\) 的 \(H^d\) 解，其余解为它加上任意核元素。
\end{theorem}
\begin{proof}
	若 \(v\in\mathcal R_P^\perp\)，则 \((Pu,v)_2=0\) 对全部光滑 \(u\) 成立，所以 \(P^*v=0\) 分布成立。正则性给出 \(v\in\mathcal N_{P^*}\)。反向由分部积分和光滑稠密性立得，故
	\(\mathcal R_P^\perp=\mathcal N_{P^*}\)。\Cref{g3:prop:elliptic-closed-range}保证值域闭；Hilbert 投影定理\cref{g3:thm:hilbert-projection}因而给出
	\(\mathcal R_P=(\mathcal R_P^\perp)^\perp\)，得到式\eqref{g3:eq:elliptic-range-orthogonal}。

	余核的具体同构把类 \([f]\) 送到 \(\Pi_{P^*}f\)：给 \(f\) 加上值域元素不改变投影，投影为零又恰好意味着 \(f\) 在值域中，故它良定义且双射。存在性随值域等式成立；减去 \(\Pi_Pu\) 给出归一化解，两个这样的解之差同时属于核和核正交补，只能为零。
\end{proof}

\begin{corollary}\label{g3:cor:elliptic-range-hk}
	设 \(M\) 为非空闭光滑黎曼流形，\(E,F\to M\) 为带固定底流形的光滑黎曼度量，以及两丛的光滑纤维度量和参考联络的同秩有限秩实或复光滑向量丛，\(P:C^\infty(E)\to C^\infty(F)\) 为 \(d\ge1\) 阶光滑椭圆算子。记 \(P^*\) 为形式伴随，\(\mathcal N_{P^*}=\ker(P^*:C^\infty(F)\to C^\infty(E))\)，对每个整数 \(k\ge0\)，有
	\[
		P(H^{k+d}(E))=\{f\in H^k(F):(f,h)_2=0\text{ 对所有 }h\in\mathcal N_{P^*}\}.
	\]
	相容右端的核正交解满足 \(\|u\|_{H^{k+d}}\le C_k\|f\|_{H^k}\)。特别地，光滑右端的解光滑。
\end{corollary}
\begin{proof}
	必要性由分部积分。相容的 \(H^k\) 右端先由上一结果得到 \(H^d\) 解，再由\cref{g3:thm:global-elliptic-hk}提升到 \(H^{k+d}\)。归一化与估计来自\cref{g3:thm:kernel-orthogonal-estimate}；光滑性可由逐阶提升或直接由分布椭圆光滑性得到。
\end{proof}

这里的余核用光滑 \(\mathcal N_{P^*}\) 表示，与 Sobolev 阶数的选择相容。不过在 \(H^k\) 中，所写条件仍是 \(L^2\) 配对，不能自行替换为 \(H^k\) 正交条件。

\subsection{可解条件和 Fredholm 指数}
\label{g3:subsec:1302:04}
\begin{definition}\label{g3:def:fredholm-operator}
	设 \(X,Y\) 为实或复 Banach 空间。若有界线性映射 \(T:X\to Y\) 的核有限维、值域闭且余核 \(Y/\operatorname{ran}T\) 有限维，则称它为\term[Fredholm suanzi]{Fredholm 算子}[Fredholm operator]，定义其\term[Fredholm zhishu]{Fredholm 指数}[Fredholm index]为
	\[
		\operatorname{ind}T=\dim\ker T-\dim(Y/\operatorname{ran}T).
	\]
\end{definition}
\symidx{\(\operatorname{ind}P\)：核维数减余核维数的 Fredholm 指数}

椭圆算子的上述实现因而是 Fredholm 的，指数为
\(\dim\mathcal N_P-\dim\mathcal N_{P^*}\)。这是一个有限整数；它不要求 \(P\) 自伴，也不要求指数为零。若 \(E=F\)、\(P=P^*\)，指数才立即为零。抽象紧算子的 Fredholm 选择及定义可对照 Brezis 第6.2节和第6章评注第1项\cite[第6.2节，第6章评注第1项]{Brezis2011}。

\begin{theorem}\label{g3:thm:elliptic-fredholm-wkp}
	设 \(M^n\) 为非空闭光滑流形，\(n\ge1\)，\(E,F\) 为同秩光滑有限秩度量丛，\(P:C^\infty(E)\to C^\infty(F)\) 为光滑椭圆 \(d\) 阶算子，\(d\ge1\)。固定底流形的光滑黎曼度量 \(g\)以及各丛的光滑参考联络，Sobolev范数用这些数据定义；积分与形式伴随采用该度量的体积密度及给定的纤维度量。记 \(\mathcal N_P,\mathcal N_{P^*}\) 为 \(P\) 及其形式伴随的光滑核，\(\Pi_P\) 为到 \(\mathcal N_P\) 的 \(L^2\) 正交投影。对整数 \(k\ge0\)、\(1<p<\infty\)，有界映射
	\[
		P:W^{k+d,p}(E)\longrightarrow W^{k,p}(F)
	\]
	为 Fredholm 算子，核为 \(\mathcal N_P\)，值域为
	\[
		\{f\in W^{k,p}(F):\textstyle\int_M\langle f,h\rangle\mu_g=0
		\text{ 对全部 }h\in\mathcal N_{P^*}\}.
	\]
	余核由 \(\mathcal N_{P^*}\) 表示，指数独立于 \(k,p\)。相容右端有唯一的 \(\Pi_Pu=0\) 解，并满足 \(\|u\|_{W^{k+d,p}}\le C_{k,p}\|Pu\|_{W^{k,p}}\)。
\end{theorem}
\begin{proof}
	先取 \(k=0\)。上一节的 \(L^p\) 核正交估计使值域闭：对收敛列 \(Pu_j\) 减去核投影后，\(u_j-\Pi_Pu_j\) 在 \(W^{d,p}\) 中为 Cauchy 列，与闭值域命题的证明完全相同。

	令 \(p'=p/(p-1)\)。由积分对偶\cref{g3:lem:lp-duality-reflexivity}及纤维度量，每个 \(L^p(F)\) 上的连续线性泛函都写成
	\(f\mapsto\int_M\langle f,h\rangle\mu_g\)，其中 \(h\in L^{p'}(F)\)。它消灭值域，当且仅当 \(P^*h=0\) 分布成立；于是 \(h\) 光滑且属于 \(\mathcal N_{P^*}\)。反之每个这样的 \(h\) 都消灭值域。若 \(f\) 不在闭值域中，Hahn--Banach 分离\cref{g3:lem:bounded-functional-extension}给出一个在值域上为零而在 \(f\) 上非零的泛函，故必有上述某个积分不为零。这就证明了 \(k=0\) 的值域等式。

	取 \(\mathcal N_{P^*}\) 的 \(L^2\) 正交归一基，积分配对构成满射 \(L^p(F)\to\mathbb K^{\dim\mathcal N_{P^*}}\)：将右端取为该光滑基的线性组合即可得到任意坐标。因此值域的余维恰为这个维数。

	对于 \(k>0\)，相容的 \(W^{k,p}\) 右端先有 \(W^{d,p}\) 解，再由\cref{g3:thm:global-elliptic-wkp}提升正则性。上述有限个积分在 \(W^{k,p}\) 中仍连续且联合满射，所以值域闭、余维相同。核光滑且不随空间改变，唯一性、归一化和估计均由上一节得到。
\end{proof}

\begin{example}\label{g3:exm:circle-fredholm-compatibility}
	在 \(S^1=\mathbb R/(2\pi\mathbb Z)\) 上，\(P=\partial_x\) 的形式伴随为 \(-\partial_x\)，两者的核都由常数生成，指数为零。方程 \(u'=f\) 有周期解当且仅当 \(\int_0^{2\pi}f=0\)；取
	\[
		u(x)=\int_0^x f(t)\,\mathrm dt-
		\frac1{2\pi}\int_0^{2\pi}\int_0^y f(t)\,\mathrm dt\,\mathrm dy
	\]
	即得唯一零均值解。这里可解性并不要求算子自伴。对 \(\Delta_g\) 而言，核在每个连通分支上由常数构成，因此相容条件是逐个分支的积分为零，不能只检查总积分。
\end{example}

\subsection{指数对连续扰动的稳定性}
\label{g3:subsec:1302:05}
核与余核的维数可能同时变化，差却保持不变。把无限维可逆部分先消去，会留下一个有限维映射；它的秩怎样变化都不改变“核维数减余维数”。

\begin{theorem}\label{g3:thm:index-local-constancy}
	设 \(M^n\) 为非空闭光滑流形，\(n\ge1\)，\(E,F\) 为同秩光滑有限秩度量丛，\(P:C^\infty(E)\to C^\infty(F)\) 为光滑椭圆 \(d\) 阶算子，\(d\ge1\)。固定底流形的光滑黎曼度量以及各丛的光滑参考联络，Sobolev范数用这些数据定义；积分与形式伴随采用该度量的体积密度及给定的纤维度量。把它看成 \(X=H^d(E)\to Y=L^2(F)\) 的有界映射。则存在 \(\delta>0\)，使每个有界线性映射 \(Q:X\to Y\) 满足 \(\|Q-P\|<\delta\) 时，\(Q\) 为 Fredholm 算子且
	\[
		\operatorname{ind}Q=\operatorname{ind}P.
	\]
	特别地，在固定 \(M,E,F\) 上，阶数为 \(d\) 的光滑椭圆算子组成的算子范数连续道路 \(P_t:H^d(E)\to L^2(F)\)、\(0\le t\le1\)，具有恒定指数。
\end{theorem}
\begin{proof}
	记 \(K=\mathcal N_P\)、\(C=\mathcal N_{P^*}\)、\(R=\mathcal R_P\)，并令 \(X_0=\ker\Pi_P\cap H^d(E)\)。我们已经证明有连续直和
	\[
		X=X_0\oplus K,\qquad Y=R\oplus C,
	\]
	以及有界可逆映射 \(A=P|_{X_0}:X_0\to R\)。按这些分解把 \(Q\) 写成
	\[
		Q=\begin{pmatrix}A_Q&B_Q\\ C_Q&D_Q\end{pmatrix}.
	\]
	投影均有界，所以 \(Q\) 充分接近 \(P\) 时，\(\|A^{-1}(A_Q-A)\|<1\)。Neumann 级数给出 \(A_Q\) 的有界逆。作有界可逆的行、列消元，得到
	\begin{equation}\label{g3:eq:fredholm-block-reduction}
		\begin{pmatrix}I&0\\-C_QA_Q^{-1}&I\end{pmatrix}
		Q
		\begin{pmatrix}I&-A_Q^{-1}B_Q\\0&I\end{pmatrix}
		=\begin{pmatrix}A_Q&0\\0&S_Q\end{pmatrix},
		\quad S_Q=D_Q-C_QA_Q^{-1}B_Q:K\to C.
	\end{equation}
	右端的第一块可逆，第二块作用在有限维空间之间，故值域为闭空间 \(R\oplus\operatorname{ran}S_Q\)，核与余核均有限维。可逆行、列变换保持这些性质和维数，因而
	\[
		\operatorname{ind}Q
		=\dim\ker S_Q-\bigl(\dim C-\operatorname{rank}S_Q\bigr)
		=\dim K-\dim C.
	\]
	这就是 \(\operatorname{ind}P\)。沿连续道路，每一点附近都可应用同一论证，故指数局部常值；区间连通使它处处相同。
\end{proof}

分块消元把所有无限维求解集中到同构 \(A_Q\)，剩下的有限维映射 \(S_Q\) 决定核与余核怎样跳跃，如图\cref{g3:fig:fredholm-reduction}。其秩同时从源维数与靶维数中扣除，因此在指数里抵消。
\begin{figure}[htbp]
	\centering
	% !TeX root = ../Book3.tex
\begin{tikzcd}[row sep=1.1em,column sep=4em]
	X_0 \arrow[r,"A_Q", "\simeq"'] & R \\
	K \arrow[r,"S_Q"'] & C
\end{tikzcd}

	\caption[可逆部分与有限维剩余方程]{可逆坐标变换后，两块独立作用。上方是闭补空间之间的同构，下方保留有限维核与余核的变化。}
	\label{g3:fig:fredholm-reduction}
\end{figure}

\begin{corollary}\label{g3:cor:lower-order-index-invariance}
	设 \(M\) 为闭光滑流形，\(E,F\) 为同秩光滑有限秩度量丛，\(P:C^\infty(E)\to C^\infty(F)\) 是光滑椭圆 \(d\) 阶算子，\(d\ge1\)。若 \(B:C^\infty(M;E)\to C^\infty(M;F)\) 为光滑系数、阶数至多 \(d-1\) 的算子，则 \(\operatorname{ind}(P+B)=\operatorname{ind}P\)。
\end{corollary}
\begin{proof}
	\(P_t=P+tB\) 的主符号与 \(P\) 相同，所以整条道路椭圆；又有 \(\|P_t-P_s\|_{H^d\to L^2}\le |t-s|\|B\|\)。应用上一结果即可。
\end{proof}

\begin{example}\label{g3:exm:kernel-jump-circle}
	在复值周期函数上取 \(P_t=-\mathrm i\partial_x+t\)，\(|t|<1/2\)。齐次解为 \(u(x)=a e^{-\mathrm itx}\)，它周期当且仅当 \(t=0\) 或 \(a=0\)。算子形式自伴，故核与余核在 \(t=0\) 同时从零维跳到一维，指数始终为零。\(t\ne0\) 时对常数右端一有 \(P_t(1/t)=1\)，因此逆算子的范数至少按 \(1/|t|\) 增长。这解释了为什么指数稳定并不保证解对参数一致有界。
\end{example}
% 纲要条目排为编号小节标题，并进入目录与书签；说明段仍保留于纲要注释中。

% !TeX root = ../../Book3.tex
% 纲要标题：### 13.3 Green 算子和谱
% 纲要位置：几何学计划纲要.md，第 2627 行。
% 正文状态：初稿；尚未正式审核。

% 纲要内容开始
% 
% * 广义逆及有限秩投影
% * 自伴非负算子的紧预解式
% * 特征值离散性和光滑特征截面
% * 谱展开的范数与正则性
% 
% 纲要内容结束

\section{Green 算子和谱}
\label{g3:sec:1303}

上一节把方程的障碍压缩到两个有限维空间：解可以加上核中的截面，右端必须正交于伴随核。固定正交补以后，就能把“存在一个解”变成一个连续的求解算子。对于形式自伴且非负的算子，这个求解过程还可与正交特征截面联系起来：椭圆估计控制导数，紧嵌入保证高频的逆作用趋于零，二者共同产生离散谱。

\subsection{广义逆及有限秩投影}
\label{g3:subsec:1303:01}
设 \(P:C^\infty(E)\to C^\infty(F)\) 是闭黎曼流形上的光滑椭圆算子。记
\[
	\mathcal N_P=\ker P,\qquad \mathcal N_{P^*}=\ker P^*,
	\qquad \Pi_P:L^2(E)\to \mathcal N_P,\quad
	\Pi_{P^*}:L^2(F)\to \mathcal N_{P^*},
\]
其中两个核都取光滑截面，并以 \(L^2\) 内积正交投影。若 \(e_1,\ldots,e_a\) 是 \(\mathcal N_P\) 的正交归一基，则
\[
	\Pi_Pu=\sum_{\nu=1}^a\langle u,e_\nu\rangle_{L^2}e_\nu.
\]
所以投影为有限秩算子，且把 \(L^2\) 截面送到光滑截面；每个 \(H^k\) 范数中的有界性都来自有限个光滑基向量的范数。基可以改变，正交投影本身不变。

\begin{theorem}\label{g3:thm:elliptic-green-general}
	设 \(M\) 为非空、紧且无边界的光滑黎曼流形，\(\dim M\ge1\)，\(E,F\to M\) 为同秩有限秩光滑复 Hermitian 丛，\(P:C^\infty(E)\to C^\infty(F)\) 为阶数 \(d\ge1\) 的光滑椭圆微分算子。存在唯一有界线性算子
	\[
		G_P:L^2(F)\longrightarrow H^d(E),
	\]
	使 \(G_Pf\perp \mathcal N_P\)、\(PG_Pf=(I-\Pi_{P^*})f\)。它满足
	\begin{equation}\label{g3:eq:green-identities}
		PG_P=I-\Pi_{P^*},\qquad
		G_PP=I-\Pi_P,\qquad
		G_P\Pi_{P^*}=0,\quad \Pi_PG_P=0,
	\end{equation}
	其中第二式作用于 \(H^d(E)\)。对每个整数 \(k\ge0\)，同一个算子有界地映射 \(H^k(F)\) 到 \(H^{k+d}(E)\)，特别地保持光滑性。作为 \(L^2(F)\to L^2(E)\) 的算子，它是紧的。若另有 \(E=F\)、\(P=P^*\)，则 \(G_P\) 在 \(L^2(E)\) 上自伴；在此条件下，若 \(\langle Pu,u\rangle_{L^2}\ge0\) 对所有光滑 \(u\) 成立，则 \(G_P\) 也非负。
\end{theorem}
\begin{proof}
	对 \(f\in L^2(F)\)，令 \(f_0=(I-\Pi_{P^*})f\)。由\cref{g3:thm:elliptic-fredholm}，\(f_0\) 属于 \(P:H^d(E)\to L^2(F)\) 的像。选择一个解 \(v\)，再取 \(u=v-\Pi_Pv\)，便得唯一的正交核解：若另有同样的解，两解之差既在 \(\mathcal N_P\) 中又正交于它，故为零。定义 \(G_Pf=u\)。唯一性保证线性，\cref{g3:thm:kernel-orthogonal-estimate}给出
	\[
		\|G_Pf\|_{H^d}\le C\|f_0\|_2\le C\|f\|_2.
	\]
	对 \(v\in H^d(E)\)，\(Pv\perp \mathcal N_{P^*}\)，而 \(v-\Pi_Pv\) 是方程 \(Pu=Pv\) 的正交核解，因此 \(G_PPv=v-\Pi_Pv\)。其余恒等式由定义直接得到。

	若 \(f\in H^k(F)\)，\cref{g3:cor:elliptic-range-hk}使 \(u=G_Pf\) 属于 \(H^{k+d}(E)\)。结合全局估计与刚得的 \(L^2\) 界，
	\[
		\|u\|_{H^{k+d}}
		\le C_k\bigl(\|f_0\|_{H^k}+\|u\|_2\bigr)
		\le C'_k\|f\|_{H^k}.
	\]
	光滑数据属于每个 \(H^k\)，故解也属于所有整数阶 Sobolev 空间，高阶嵌入\cref{g3:cor:higher-sobolev-embedding}遂给出光滑性。由\cref{g3:thm:compact-bundle-rellich}，\(H^d(E)\Subset L^2(E)\)，所以 \(G_P\) 的 \(L^2\) 作用为紧算子。

	最后设 \(P=P^*\)。取 \(u=G_Pf\)、\(v=G_Pg\)。光滑稠密性把形式伴随恒等式延伸到 \(H^d\)；核的正交性于是给出
	\[
		\langle G_Pf,g\rangle
		=\langle u,Pv\rangle
		=\langle Pu,v\rangle
		=\langle f,G_Pg\rangle.
	\]
	若 \(P\) 非负，同一稠密性使 \(\langle Pu,u\rangle\ge0\) 对 \(u\in H^d\) 仍成立，因而
	\(\langle G_Pf,f\rangle=\langle u,Pu\rangle\ge0\)。
\end{proof}

\(G_P\) 称为所固定 \(L^2\) 度量下的\term[Green suanzi]{Green 算子}[Green operator]，也是 \(P\) 的规范广义逆。\footnote{George Green（1793--1841），英国数学家，研究势论与数学物理，Green 函数及有关积分公式以其命名。}它在可解右端上选择唯一的最小 \(L^2\) 范数解，因为其余解为 \(G_Pf+h\)、\(h\in \mathcal N_P\)，且
\(\|G_Pf+h\|_2^2=\|G_Pf\|_2^2+\|h\|_2^2\)。在不可解右端上，算子先去掉伴随核分量；式\eqref{g3:eq:green-identities}准确记录了这个有限维误差。

\subsection{自伴非负算子的紧预解式}
\label{g3:subsec:1303:02}
以下取 \(E=F\)，并假设 \(P=P^*\) 且 \(P\) 非负。必须先明确 \(P\) 在 \(L^2\) 中的定义域。微分算子会损失导数，通常不是处处定义的有界 \(L^2\) 算子；将它看作
\[
	A:\mathcal D(A)=H^d(E)\subset L^2(E)\longrightarrow L^2(E)
\]
才与此前的椭圆估计相符。对于稠密定义算子，伴随域由配对决定：\(v\in\mathcal D(A^*)\) 表示存在 \(w\in L^2\) 使 \(\langle Au,v\rangle=\langle u,w\rangle\) 对所有 \(u\in\mathcal D(A)\) 成立，此时定义 \(A^*v=w\)。自伴要求域与作用同时相等。

\begin{theorem}\label{g3:thm:positive-elliptic-realization}
	设 \(M\) 为非空、紧且无边界的光滑黎曼流形，\(\dim M\ge1\)，\(E\to M\) 为有限秩光滑复 Hermitian 丛，\(P\) 为阶数 \(d\ge1\) 的光滑椭圆算子，满足 \(P=P^*\) 及 \(\langle Pu,u\rangle\ge0\) 对所有 \(u\in C^\infty(E)\) 成立。令 \(A\) 为 \(P\) 在 \(L^2(E)\) 中定义域为 \(H^d(E)\) 的实现。则 \(A\) 是稠密定义、闭且非负的自伴算子，并且
	\[
		R=(A+I)^{-1}:L^2(E)\longrightarrow H^d(E)
	\]
	存在且有界。作为 \(L^2\) 算子，\(R\) 紧、自伴、单射，满足 \(\|R\|_{L^2\to L^2}\le1\)，并且 \(\langle Rf,f\rangle>0\) 对每个非零 \(f\) 成立。
\end{theorem}
\begin{proof}
	\Cref{g3:prop:elliptic-adjoint-domain}已经证明：\(A\) 稠定且闭，其 Hilbert 伴随的域为 \(H^d(E)\)，作用为形式伴随 \(P^*\)。这里 \(P^*=P\)，故 \(A^*=A\)，包括域也相同。光滑截面在 \(H^d\) 中稠密，\(P:H^d\to L^2\) 有界，所以光滑截面上的非负性取极限后对整个定义域成立。

	算子 \(P+I\) 仍椭圆、形式自伴，而且
	\[
		\langle(P+I)u,u\rangle\ge\|u\|_2^2.
	\]
	故它的核为零。\Cref{g3:thm:elliptic-fredholm}把像的正交补识别为同一个零核，因而 \(P+I:H^d\to L^2\) 为双射。令 \(u=Rf\)，配对方程给出 \(\|u\|_2\le\|f\|_2\)，再由全局椭圆估计得到 \(\|u\|_{H^d}\le C\|f\|_2\)。紧性随即由 \(H^d\Subset L^2\) 得到。

	设 \(u=Rf\)、\(v=Rg\)。则
	\[
		\langle Rf,g\rangle
		=\langle u,(A+I)v\rangle
		=\langle(A+I)u,v\rangle
		=\langle f,Rg\rangle.
	\]
	这证明 \(R\) 自伴。又有
	\(\langle Rf,f\rangle=\langle u,(A+I)u\rangle\ge\|u\|_2^2\)，非零 \(f\) 对应非零 \(u\)，故得到严格正性和单射性。
\end{proof}

若一个复数 \(z\) 使 \(A-zI:\mathcal D(A)\to L^2\) 为双射，且其逆在 \(L^2\) 上有界，则称 \(z\) 属于\term[yujieji]{预解集}[resolvent set] \(\rho(A)\)，逆 \((A-zI)^{-1}\) 称为\term[yujieshi]{预解式}[resolvent]。补集 \(\sigma(A)=\mathbb C\setminus\rho(A)\) 称为\term[pu]{谱}[spectrum]。刚才的定理给出 \(-1\in\rho(A)\)，并证明这个预解式紧。加上 \(I\) 的作用是统一消去可能存在的零特征值，随后再从 \(R\) 的特征值反求 \(A\) 的特征值。

\subsection{特征值离散性和光滑特征截面}
\label{g3:subsec:1303:03}
先把紧性怎样产生特征向量这一泛函分析步骤完整写出。下面只需要非负自伴情形；零空间可以无限维，而每个正特征值的空间必为有限维。

\begin{theorem}[紧非负自伴算子的谱分解]\label{g3:thm:compact-positive-spectrum}
	设 \(H\) 为可分的实或复 Hilbert 空间，\(T:H\to H\) 为紧自伴算子，且 \(\langle Tu,u\rangle\ge0\)。则 \(H\) 有由 \(T\) 的特征向量组成的正交归一基。每个正特征值具有有限重数；按重数排列的正特征值若有无限多个，只能趋于零。特别地，若 \(H\) 无限维且 \(\ker T=0\)，则存在正交归一基 \((e_j)_{j\ge1}\) 及
	\[
		\mu_1\ge\mu_2\ge\cdots>0,\qquad
		\mu_j\longrightarrow0,\qquad Te_j=\mu_je_j.
	\]
\end{theorem}
\begin{proof}
	证明的核心是让单位球上二次型的上确界达到。先记
	\(a=\sup_{\|u\|=1}\langle Tu,u\rangle\)。非负 Hermitian 型 \(b(u,v)=\langle Tu,v\rangle\) 满足
	\[
		|b(u,v)|^2\le b(u,u)b(v,v).
	\]
	确实，将 \(b(u-tv,u-tv)\ge0\) 对标量 \(t\) 配方即可；当 \(b(v,v)=0\) 时，让 \(t\) 沿交叉项的相位取任意大模长，必有 \(b(u,v)=0\)。所以
	\(|\langle Tu,v\rangle|\le a\|u\|\|v\|\)，即 \(\|T\|\le a\)；反向不等式直接成立，故 \(a=\|T\|\)。若 \(T\ne0\)，则 \(a>0\)。

	取单位向量列 \(u_\ell\) 使 \(\langle Tu_\ell,u_\ell\rangle\to a\)。由\cref{g3:lem:hilbert-weak-minimization-tools}，抽子列使 \(u_\ell\rightharpoonup u\)，\(\|u\|\le1\)；紧性再使 \(Tu_\ell\) 强收敛。它的极限必为 \(Tu\)，因为对任意 \(v\)，\(\langle Tu_\ell,v\rangle=\langle u_\ell,Tv\rangle\to\langle Tu,v\rangle\)。于是
	\[
		a=\langle Tu,u\rangle\le a\|u\|^2\le a.
	\]
	因此 \(\|u\|=1\)，并且 \(\langle(aI-T)u,u\rangle=0\)。算子 \(aI-T\) 非负，对它使用刚证的半内积不等式，得 \(\langle(aI-T)u,v\rangle=0\) 对所有 \(v\) 成立。因此 \(Tu=au\)。

	取 \(e_1=u\)。自伴性保证 \(e_1^\perp\) 在 \(T\) 下不变：\(\langle Tv,e_1\rangle=\langle v,Te_1\rangle=0\)。在这个闭子空间上重复最大化过程；只要限制算子非零，就得到新的单位特征向量 \(e_j\)，其特征值 \(\mu_j\) 为剩余空间上 \(T\) 的范数。因此 \(\mu_j\) 非增。若有无限多个 \(\mu_j\ge\varepsilon>0\)，则
	\[
		\|Te_i-Te_j\|^2=\mu_i^2+\mu_j^2\ge2\varepsilon^2
		\quad(i\ne j),
	\]
	与紧性矛盾。这同时说明正特征值的重数有限，以及无限序列 \(\mu_j\to0\)。

	令 \(V\) 为所得正特征向量的闭线性张成空间。在 \(V^\perp\) 上若 \(T\) 非零，刚证的最大化论证又产生正特征值 \(b>0\)。无限过程的每个剩余空间都包含 \(V^\perp\)，因而每个 \(\mu_j\ge b\)，与趋零矛盾；有限过程中则与停止时限制算子为零矛盾。因此 \(T|_{V^\perp}=0\)。另一方面，\(\ker T\) 正交于所有正特征向量，故 \(V^\perp=\ker T\)。在可分闭子空间 \(\ker T\) 的可数稠密集上施行逐次正交化，丢掉为零的余向量，即得它的正交归一基；与 \((e_j)\) 合在一起得到所需基。

	最后，任意 \(x\in H\) 到前 \(N\) 个基向量张成空间的正交投影为 \(\sum_{j\le N}\langle x,e_j\rangle e_j\)。其误差与该空间正交，平方范数为 \(\|x\|^2-\sum_{j\le N}|\langle x,e_j\rangle|^2\)。基的闭张成性使误差趋零，故展开在范数中收敛且满足平方和恒等式。这也说明上述“基”具有实际的级数展开含义。
\end{proof}

紧自伴分解可对照 Brezis 第6.4节的定理6.11。\cite[第6.4节，定理6.11]{Brezis2011}这里通过弱收敛与紧性直接取得二次型最大值，使所需的正算子版本不依赖一般有界算子的谱理论。

\begin{theorem}\label{g3:thm:elliptic-discrete-spectrum}
	设 \(M\) 为非空、紧且无边界的光滑黎曼流形，\(\dim M\ge1\)，\(E\to M\) 为正秩光滑复 Hermitian 丛。设 \(P\) 为光滑、形式自伴、非负、椭圆的 \(d\ge1\) 阶微分算子，令 \(A\) 为其定义域 \(H^d(E)\) 的 \(L^2\) 实现。则存在光滑截面组成的 \(L^2\) 正交归一基 \((e_j)_{j\ge1}\)，使
	\[
		Pe_j=\lambda_je_j,\qquad
		0\le\lambda_1\le\lambda_2\le\cdots,
		\qquad \lambda_j\longrightarrow+\infty.
	\]
	各特征值均具有有限重数，\(\sigma(A)=\{\lambda_j:j\ge1\}\)。对 \(u=\sum_jc_je_j\in L^2(E)\)，有
	\begin{equation}\label{g3:eq:spectral-domain}
		u\in H^d(E)
		\quad\Longleftrightarrow\quad
		\sum_j(1+\lambda_j)^2|c_j|^2<\infty;
	\end{equation}
	此时 \(Au=\sum_j\lambda_jc_je_j\)，级数在 \(L^2\) 中收敛。
\end{theorem}
\begin{proof}
	\(L^2(E)\) 可分：有限坐标平凡化中，以有理阶梯函数逼近各分量，再用固定单位分解拼合，可得可数稠密集；它又因正维坐标球中可选择无限多个线性独立紧支截面而无限维。将紧谱定理用于\cref{g3:thm:positive-elliptic-realization}的 \(R=(A+I)^{-1}\)。因 \(R\) 单射，所得基向量的特征值 \(\mu_j\) 均为正，\(\mu_j\to0\)。从 \(e_j=\mu_j^{-1}Re_j\) 知 \(e_j\in H^d\)，且
	\[
		Ae_j=(\mu_j^{-1}-1)e_j.
	\]
	令 \(\lambda_j=\mu_j^{-1}-1\)。由非负性，\(\lambda_j\ge0\)，其有限重数与趋于正无穷均由 \(\mu_j\) 的性质得到。方程 \(Pe_j=\lambda_je_j\) 反复应用\cref{g3:thm:elliptic-distribution-l2}，使 \(e_j\) 依次属于 \(H^{2d},H^{3d},\ldots\)，故由高阶嵌入\cref{g3:cor:higher-sobolev-embedding}有光滑代表。

	若 \(u\in H^d\)，形式伴随恒等式给出 \(\langle Au,e_j\rangle=\lambda_jc_j\)。平方和恒等式说明 \(\sum_j\lambda_j^2|c_j|^2=\|Au\|_2^2\)，再加上 \(u\in L^2\)，即得式\eqref{g3:eq:spectral-domain}的必要性。反之，若该带权平方和有限，则部分和 \(u_N=\sum_{j\le N}c_je_j\) 及 \(Pu_N\) 分别在 \(L^2\) 中为 Cauchy 列。\Cref{g3:thm:global-elliptic-hk}给出
	\[
		\|u_N-u_K\|_{H^d}
		\le C\bigl(\|P(u_N-u_K)\|_2+\|u_N-u_K\|_2\bigr).
	\]
	故它们在 \(H^d\) 中收敛；其 \(L^2\) 极限是 \(u\)，这就证明充分性及 \(Au\) 的展开。

	每个 \(\lambda_j\) 显然不属于预解集。若 \(z\) 不等于任何 \(\lambda_j\)，由 \(\lambda_j\to+\infty\)，有
	\[
		\inf_j|\lambda_j-z|>0,\qquad
		\sup_j\frac{1+\lambda_j}{|\lambda_j-z|}<\infty.
	\]
	因此对 \(f=\sum_jf_je_j\)，级数
	\[
		u=\sum_j\frac{f_j}{\lambda_j-z}e_j
	\]
	定义一个 \(H^d\) 截面，满足 \((A-zI)u=f\)，且 \(L^2\) 范数由 \(C_z\|f\|_2\) 控制。逐个系数又给出唯一性，故 \(z\in\rho(A)\)。这证明了整个谱的描述，而不只是列出若干特征值。
\end{proof}

\begin{example}\label{g3:exm:circle-elliptic-spectrum}
	在周长 \(2\pi\) 的圆周上，取 \(P=-\partial_x^2\)。解常微分方程 \(-u''=\lambda u\) 并施加函数及一阶导数的周期条件，得到零特征空间为常数，正特征值为 \(k^2\)、\(k\ge1\)，每个复特征空间由 \(e^{\mathrm ikx},e^{-\mathrm ikx}\) 张成。上述定理保证除此之外没有缺失的 \(L^2\) 方向。因此
	\[
		e_k(x)=(2\pi)^{-1/2}e^{\mathrm ikx},\qquad k\in\mathbb Z,
	\]
	构成正交归一基，\(Re_k=(1+k^2)^{-1}e_k\)，而 \(G_Pe_0=0\)、\(G_Pe_k=k^{-2}e_k\) 对 \(k\ne0\) 成立。原算子在高频上增长，两个逆算子的高频作用却衰减，这正是紧性的具体表现。
\end{example}

\subsection{谱展开的范数与正则性}
\label{g3:subsec:1303:04}
特征展开不仅描述 \(L^2\) 函数，还可记录导数范数。最直接的版本按算子的整数次幂计数；它已经足以识别光滑截面，并控制谱截断在每个固定导数阶数中的误差。

\begin{theorem}\label{g3:thm:elliptic-spectral-sobolev}
	设 \(M\) 为非空、紧且无边界的正维光滑黎曼流形，\(E\to M\) 为正秩光滑复 Hermitian 丛，\(P\) 为光滑、形式自伴、非负的 \(d\ge1\) 阶椭圆微分算子。取\cref{g3:thm:elliptic-discrete-spectrum}的特征基，写 \(u=\sum_jc_je_j\in L^2(E)\)。对每个整数 \(q\ge0\)，有
	\begin{equation}\label{g3:eq:spectral-sobolev-weight}
		u\in H^{qd}(E)
		\quad\Longleftrightarrow\quad
		\sum_j(1+\lambda_j)^{2q}|c_j|^2<\infty,
	\end{equation}
	并且存在只依赖 \(P,q\) 及所固定几何数据的正数 \(c_q,C_q\)，使
	\[
		c_q\|u\|_{H^{qd}}^2
		\le\sum_j(1+\lambda_j)^{2q}|c_j|^2
		\le C_q\|u\|_{H^{qd}}^2.
	\]
	特别地，\(u\) 有光滑代表当且仅当这些带权平方和对每个 \(q\) 都有限；此时特征展开在每个 \(C^r\) 范数中收敛。
\end{theorem}
\begin{proof}
	\(q=0\) 正是正交归一展开的平方和恒等式。设 \(q\ge1\)，令 \(B_q=(P+I)^q\)，这是阶数 \(qd\) 的光滑椭圆算子：其最高阶主符号为 \(\sigma_d(P)^q\)，在非零余切方向上仍可逆。对 \(u\in H^{qd}\)，光滑稠密性及逐次形式伴随给出
	\[
		\langle B_qu,e_j\rangle=(1+\lambda_j)^qc_j,
		\qquad
		\|B_qu\|_2^2=\sum_j(1+\lambda_j)^{2q}|c_j|^2.
	\]
	微分算子的有界性给出右侧至多为 \(C_q\|u\|_{H^{qd}}^2\)。另一方面，对 \(B_q\) 使用\cref{g3:thm:global-elliptic-hk}，得
	\[
		\|u\|_{H^{qd}}\le C\bigl(\|B_qu\|_2+\|u\|_2\bigr)
		\le2C\|B_qu\|_2,
	\]
	最后一步使用 \((1+\lambda_j)^q\ge1\)。这证明了范数等价及必要性。

	若带权平方和有限，对光滑有限和 \(u_N=\sum_{j\le N}c_je_j\) 应用刚才的估计，便知 \(u_N\) 在 \(H^{qd}\) 中为 Cauchy 列。其极限的 \(L^2\) 展开为 \(u\)，故 \(u\in H^{qd}\)，并在取极限后保留全部恒等式与不等式。

	光滑截面在紧流形上具有所有整数阶 Sobolev 范数。反之，若每个 \(q\) 的带权平方和有限，对任意整数 \(r\ge0\)，选择 \(qd>r+\dim M/2\)。高阶 Sobolev 嵌入\cref{g3:cor:higher-sobolev-embedding}给出 \(H^{qd}(E)\hookrightarrow C^r(E)\)。因此 \(u\) 的各个连续代表相容，并具有任意阶连续导数；同时 \(u_N\to u\) 在 \(H^{qd}\) 中收敛，故也在 \(C^r\) 中收敛。
\end{proof}

上述正则性判据中的衰减是带权平方可和条件，导数阶数由权重的幂明确计量。它还使 Green 算子具有可直接使用的谱公式：
\begin{equation}\label{g3:eq:green-spectral-series}
	G_Pf=\sum_{\lambda_j>0}\frac{\langle f,e_j\rangle}{\lambda_j}e_j,
	\qquad
	\Pi_Pf=\sum_{\lambda_j=0}\langle f,e_j\rangle e_j.
\end{equation}
对 \(f\in L^2\)，第一式在 \(H^d\) 中收敛，因为正特征值有正的最小值，且 \((1+\lambda_j)/\lambda_j\) 在正谱上一致有界。所得截面满足正交核条件和方程，故由唯一性就是先前构造的 \(G_Pf\)。这把可解性、规范解与特征展开联系在同一算子上。

\begin{corollary}\label{g3:cor:spectral-truncation-error}
	设 \(M\) 为非空、紧且无边界的正维光滑黎曼流形，\(E\to M\) 为正秩光滑复 Hermitian 丛，\(P\) 为光滑、形式自伴、非负的 \(d\ge1\) 阶椭圆算子。令 \(\Pi_{\le L}\) 为其特征值不超过 \(L\ge0\) 的有限维特征空间之和上的 \(L^2\) 正交投影。对每个整数 \(q\ge0\) 及 \(u\in H^{(q+1)d}(E)\)，有
	\[
		\|(I-\Pi_{\le L})u\|_{H^{qd}}
		\le\frac{C_q}{1+L}\|u\|_{H^{(q+1)d}}.
	\]
\end{corollary}
\begin{proof}
	在 \(\lambda_j>L\) 的项上有
	\((1+\lambda_j)^{2q}\le(1+L)^{-2}(1+\lambda_j)^{2(q+1)}\)。乘以 \(|c_j|^2\) 求和，再对两端使用式\eqref{g3:eq:spectral-sobolev-weight}的范数等价，开平方即得结论。
\end{proof}
% 纲要条目排为编号小节标题，并进入目录与书签；说明段仍保留于纲要注释中。

% !TeX root = ../../Book3.tex
% 纲要标题：### 13.4 椭圆复形
% 纲要位置：几何学计划纲要.md，第 2634 行。
% 正文状态：初稿；尚未正式审核。

% 纲要内容开始
% 
% * 符号序列的正合性
% * 相应 Laplace 算子的椭圆性
% * 正合、余正合与调和分解
% * 有限维上同调及代表元
% * de Rham 和 Dolbeault 的共同机制
% 
% 纲要内容结束

\section{椭圆复形}
\label{g3:sec:1304}

外微分满足 \(\mathrm d^2=0\)，于是正合形式一定闭合，而闭合形式未必正合。这里的障碍原本定义为一个无限维空间的商。怎样用有限个光滑对象描述它？单独的外微分通常不是方形椭圆算子；把相邻两个方向的微分及其伴随合在一起，却能得到椭圆的二阶算子。上一节的 Green 算子将使每个截面分成三个正交部分，并使上同调类有唯一的调和代表。

\subsection{符号序列的正合性}
\label{g3:subsec:1304:01}
先把“相邻两次微分为零”与“每个非零频率都没有剩余障碍”分开。前者是微分算子的恒等式，后者是各纤维上的有限维线性代数条件。

\begin{definition}\label{g3:def:first-order-elliptic-complex}
	设 \(M\) 为光滑流形，\(E^0,\ldots,E^N\) 为有限秩复向量丛，\(D_q:C^\infty(M;E^q)\to C^\infty(M;E^{q+1})\) 为阶数至多一的光滑系数微分算子。若 \(D_{q+1}D_q=0\)，则
	\[
		0\longrightarrow C^\infty(E^0)
		\xrightarrow{D_0}C^\infty(E^1)
		\xrightarrow{D_1}\cdots
		\xrightarrow{D_{N-1}}C^\infty(E^N)
		\longrightarrow0
	\]
	称为一阶\term[weifen fuxing]{微分复形}[differential complex]。记 \(a_q(x,\xi)=\sigma_1(D_q)(x,\xi)\)。若对每个 \(x\in M\) 和每个非零实余切向量 \(\xi\)，序列
	\begin{equation}\label{g3:eq:elliptic-symbol-sequence}
		0\longrightarrow E_x^0\xrightarrow{a_0(x,\xi)}E_x^1
		\longrightarrow\cdots\xrightarrow{a_{N-1}(x,\xi)}E_x^N
		\longrightarrow0
	\end{equation}
	在每一项正合，就称该复形为\term[tuoyuan fuxing]{椭圆复形}[elliptic complex]。端点处约定 \(E^{-1}=E^{N+1}=0\)、\(D_{-1}=D_N=0\)。
\end{definition}

由主符号的复合公式\cref{g3:prop:symbol-calculus}，\(D_{q+1}D_q=0\) 蕴含 \(a_{q+1}a_q=0\)，所以像总包含在下一个核中。椭圆性进一步要求这两个子空间相等。式\eqref{g3:eq:elliptic-symbol-sequence}首项的正合要求 \(a_0\) 单射，末项要求 \(a_{N-1}\) 满射；不能只核对中间项。只有两个非零丛时，这个定义恰好要求唯一的主符号可逆，因而回到普通椭圆算子。

\begin{example}\label{g3:exm:gradient-curl-symbol-complex}
	在 \(\mathbb R^2\) 上，取
	\[
		D_0f=(\partial_1f,\partial_2f),\qquad
		D_1(v_1,v_2)=\partial_1v_2-\partial_2v_1.
	\]
	混合偏导可交换，故 \(D_1D_0=0\)。沿用 \(\mathrm D_j=-\mathrm i\partial_j\) 的约定，两个符号为
	\[
		a_0(\xi)c=\mathrm i c(\xi_1,\xi_2),\qquad
		a_1(\xi)v=\mathrm i(\xi_1v_2-\xi_2v_1).
	\]
	若 \(\xi\ne0\)，第一映射单射，第二映射满射，而第二映射的核正是 \(\mathbb C\xi\)。因此它是椭圆复形。若删掉最后的旋度映射，\(0\to\mathbb C\xrightarrow{a_0}\mathbb C^2\to0\) 在末项便不正合：梯度的符号只控制沿 \(\xi\) 的方向。
\end{example}

\begin{example}\label{g3:exm:nonelliptic-complex-infinite-cohomology}
	在 \(\mathbb T^2=(\mathbb R/2\pi\mathbb Z)^2\) 上，考虑两项复形 \(C^\infty(\mathbb T^2)\xrightarrow{\partial_x}C^\infty(\mathbb T^2)\)。它在 \(\xi_x=0\) 的非零方向不椭圆。核含有所有 \(f(y)\)，因而无限维。靶空间对像的商也同构于 \(C^\infty(S^1)\)：映射是取 \(x\) 方向平均，零平均函数 \(g\) 的原函数 \(\int_0^x g(t,y)\,\mathrm dt\) 则是周期光滑函数。复形恒等式本身没有带来有限维性。
\end{example}

\subsection{相应 Laplace 算子的椭圆性}
\label{g3:subsec:1304:02}
固定各丛的光滑 Hermitian 度量与底流形的光滑正密度，形式伴随由积分配对定义。对 \(E^q\) 中的一个截面，有两个可测量的方向：向前的 \(D_qu\)，以及与前一箭头对应的 \(D_{q-1}^*u\)。把两者的平方能量相加，正好产生所需的方形算子。

\ProseParagraphBreak

正密度的选择也能纳入上一节的黎曼框架。若 \(\dim M=n\ge1\)，取任意底度量 \(g_0\)，写 \(\mu=\rho\mu_{g_0}\)，其中 \(\rho>0\) 光滑；令 \(g=\rho^{2/n}g_0\)，便有 \(\mu_g=\mu\)。在紧流形上，更换光滑底度量只改变 Sobolev 范数的等价常数，因此可以直接应用上一节以黎曼体积定义的 Green 算子理论。

\begin{lemma}\label{g3:lem:complex-symbol-laplacian}
	设 \(V_-,V,V_+\) 为有限维 Hermitian 空间，\(A:V_-\to V\)、\(B:V\to V_+\) 满足 \(BA=0\)。则
	\[
		\ker(AA^*+B^*B)=\ker A^*\cap\ker B.
	\]
	特别地，\(\operatorname{im}A=\ker B\) 当且仅当 \(AA^*+B^*B\) 正定。
\end{lemma}
\begin{proof}
	对 \(v\in V\) 与自身配对，有
	\[
		\langle(AA^*+B^*B)v,v\rangle=|A^*v|^2+|Bv|^2.
	\]
	两项非负，故左侧为零恰好意味着右侧两向量都为零。若中间正合，这样的 \(v\) 属于 \(\operatorname{im}A\)，同时正交于 \(\operatorname{im}A\)，所以 \(v=0\)。反过来，由 \(BA=0\)，\(\operatorname{im}A\subset\ker B\)。任意 \(v\in\ker B\) 可分解为在 \(\operatorname{im}A\) 上的投影与正交余项；余项仍在 \(\ker B\)，又在 \(\ker A^*\)，正定性迫使余项为零。
\end{proof}

\begin{proposition}\label{g3:prop:elliptic-complex-laplacians}
	设 \(M\) 为维数 \(n\ge1\) 的紧致无边界光滑流形，\((E^\bullet,D)\) 为其上的一阶椭圆复形，各丛带光滑 Hermitian 度量，积分采用固定光滑正密度。则
	\begin{equation}\label{g3:eq:complex-laplacian}
		\Delta_q=D_{q-1}D_{q-1}^*+D_q^*D_q
	\end{equation}
	是 \(E^q\) 上形式自伴、非负的二阶椭圆算子。相对于任意固定底度量，其主符号满足
	\[
		\sigma_2(\Delta_q)=a_{q-1}a_{q-1}^*+a_q^*a_q,
		\qquad
		\langle\sigma_2(\Delta_q)(x,\xi)v,v\rangle
		\ge c_q|\xi|^2|v|^2
	\]
	对某个 \(c_q>0\) 成立。其核是有限维光滑空间
	\[
		\mathcal H_D^q
		=\{u\in C^\infty(E^q):D_qu=0,\ D_{q-1}^*u=0\}.
	\]
\end{proposition}
\begin{proof}
	符号公式来自复合及伴随的符号运算，上一引理与符号正合性保证它对每个 \(\xi\ne0\) 正定。在单位余切球面与单位纤维球面上，连续的正函数取得正的最小值；这些球面丛在紧底空间上紧。再用符号的二次齐次性即得所述统一下界。

	形式伴随的定义与无边界条件给出，对光滑 \(u\)，
	\begin{equation}\label{g3:eq:complex-energy-identity}
		(\Delta_qu,u)_{L^2}
		=\|D_qu\|_2^2+\|D_{q-1}^*u\|_2^2.
	\end{equation}
	每个乘积项均形式自伴，因此 \(\Delta_q\) 形式自伴且非负。上一节的\cref{g3:thm:elliptic-green-general}适用于这个方形椭圆算子，给出有限维光滑核；式\eqref{g3:eq:complex-energy-identity}又把该核识别为同时满足两个一阶方程的截面。反向包含直接代入式\eqref{g3:eq:complex-laplacian}即可。
\end{proof}

\(\mathcal H_D^q\) 中的元素称为该复形的\term[tiaohe jiemian]{调和截面}[harmonic section]。这一说法同时包含闭合与余闭合两个条件。定义形式伴随使用正密度即可，以上论证没有要求流形可定向。

\subsection{正合、余正合与调和分解}
\label{g3:subsec:1304:03}
令 \(H_q\) 为到 \(\mathcal H_D^q\) 的 \(L^2\) 正交投影，\(G_q\) 为 \(\Delta_q\) 的 Green 算子。由\cref{g3:thm:elliptic-green-general}，二者保持光滑截面，且
\[
	\Delta_qG_q=G_q\Delta_q=I-H_q,\qquad
	G_q:H^k(E^q)\longrightarrow H^{k+2}(E^q),\quad k\ge0.
\]
这里第一式按各自定义域理解；作用于光滑截面时全部运算均有意义。要使这个解算子反映复形结构，须先证明它与相邻微分相容。

\begin{lemma}\label{g3:lem:complex-green-commutation}
	设 \(M\) 为维数 \(n\ge1\) 的紧致无边界光滑流形，\((E^\bullet,D)\) 为带光滑 Hermitian 度量与正密度的一阶椭圆复形。其 Laplace 算子、调和投影与 Green 算子在光滑截面上满足
	\begin{align*}
		\Delta_{q+1}D_q&=D_q\Delta_q,
		&\Delta_qD_q^*&=D_q^*\Delta_{q+1},\\
		H_{q+1}D_q&=0,
		&D_qH_q&=0,\\
		D_qG_q&=G_{q+1}D_q,
		&D_q^*G_{q+1}&=G_qD_q^*.
	\end{align*}
\end{lemma}
\begin{proof}
	展开 \(\Delta_{q+1}D_q\) 与 \(D_q\Delta_q\)，两者都剩下 \(D_qD_q^*D_q\)，因为 \(D_{q+1}D_q=D_qD_{q-1}=0\)。取形式伴随得到第二个交换关系。调和截面被 \(D_q\) 消去，所以 \(D_qH_q=0\)。若 \(h\in\mathcal H_D^{q+1}\)，则
	\((D_qu,h)=(u,D_q^*h)=0\)，故 \(H_{q+1}D_q=0\)。

	现在比较 \(D_qG_qu\) 与 \(G_{q+1}D_qu\)。第一交换关系给出
	\[
		\Delta_{q+1}D_qG_qu=D_q(u-H_qu)=D_qu.
	\]
	又由 \(H_{q+1}D_q=0\)，Green 恒等式给出
	\(\Delta_{q+1}G_{q+1}D_qu=D_qu\)。两者都正交于 \(\mathcal H_D^{q+1}\)：第一项因它在 \(D_q\) 的像中，第二项因 Green 算子的取值约定。正交核后的解唯一，因此两者相等。对 \(D_q^*\) 重复相同论证，或对刚证的等式使用 \(G_q^*=G_q\) 并与任意光滑截面配对，即得最后一个关系。
\end{proof}

\term[Hodge fenjie]{Hodge 分解}[Hodge decomposition]\footnote{William Vallance Douglas Hodge（1903--1975），苏格兰数学家，研究代数几何、微分几何与调和积分；用调和形式表示上同调类的理论以其命名。}把 Green 恒等式中的二阶项拆成两个一阶像。下述直和的正交性是 \(L^2\) 正交性；光滑截面空间本身没有被当成完备 Hilbert 空间。

\begin{theorem}\label{g3:thm:elliptic-complex-hodge-decomposition}
	设 \(M\) 为维数 \(n\ge1\) 的紧致无边界光滑流形，\((E^\bullet,D)\) 为带光滑 Hermitian 度量与正密度的一阶椭圆复形。则有 \(L^2\) 正交直和
	\begin{equation}\label{g3:eq:smooth-complex-hodge}
		C^\infty(E^q)=\mathcal H_D^q
		\oplus D_{q-1}C^\infty(E^{q-1})
		\oplus D_q^*C^\infty(E^{q+1}),
	\end{equation}
	其中任意 \(u\) 的三个分量分别为
	\begin{equation}\label{g3:eq:explicit-complex-hodge}
		u=H_qu+D_{q-1}(D_{q-1}^*G_qu)+D_q^*(D_qG_qu).
	\end{equation}
	还成立
	\[
		L^2(E^q)=\mathcal H_D^q
		\oplus D_{q-1}H^1(E^{q-1})
		\oplus D_q^*H^1(E^{q+1}),
	\]
	右侧三个子空间均闭。此时式\eqref{g3:eq:explicit-complex-hodge}仍成立，括号内的两个原像属于 \(H^1\)。
\end{theorem}
\begin{proof}
	在光滑截面上，\(I-H_q=\Delta_qG_q\) 直接给出式\eqref{g3:eq:explicit-complex-hodge}。两个像之间正交，因为
	\[
		(D_{q-1}v,D_q^*w)=(D_qD_{q-1}v,w)=0.
	\]
	它们又都正交于调和空间，因为调和截面被相应伴随或微分消去。因此每个 \(u\) 的三个分量唯一。原像本身可以加上相应算子的核，故唯一的是三个分量，而非所有表示这些分量的原像。

	记 \(\mathsf E_q=D_{q-1}D_{q-1}^*G_q\)、\(\mathsf C_q=D_q^*D_qG_q\)。由于 \(G_q:L^2\to H^2\) 有界，两个一阶算子依次作用后，这两映射在 \(L^2\) 上有界。光滑情形的正交分解表明 \(H_q,\mathsf E_q,\mathsf C_q\) 在稠密光滑子空间上是互相正交的投影，和为恒等算子。取 \(L^2\) 极限后，同样的恒等式与正交性在全空间成立，所以三个像都是闭子空间。

	还须识别像的实际定义域。\(\mathsf E_qu=D_{q-1}(D_{q-1}^*G_qu)\)，括号内属于 \(H^1\)，故 \(\operatorname{im}\mathsf E_q\subset D_{q-1}H^1\)。反过来，若 \(v\in H^1(E^{q-1})\)，取光滑 \(v_j\to v\) 于 \(H^1\)。光滑分解给出 \(\mathsf E_qD_{q-1}v_j=D_{q-1}v_j\)，由 \(D_{q-1}:H^1\to L^2\) 与 \(\mathsf E_q:L^2\to L^2\) 的连续性得 \(\mathsf E_qD_{q-1}v=D_{q-1}v\)。因此两像相等。\(\operatorname{im}\mathsf C_q=D_q^*H^1\) 的证明相同，结论全部成立。
\end{proof}

两个微分像分别称为正合部分与\term[yuzhenghe]{余正合}[coexact]部分。式\eqref{g3:eq:explicit-complex-hodge}既给出存在性，也给出了可选择的原像；其中 Green 算子的两阶增益，保证原像比 \(u\) 多一阶正则性。

\subsection{有限维上同调及代表元}
\label{g3:subsec:1304:04}
在代数定义中，上同调只记住闭截面对正合截面的等价类。调和分解将从每个类中选出一个同时满足椭圆方程的截面，使这个商有有限维的具体模型。

\begin{definition}\label{g3:def:elliptic-complex-cohomology}
	设 \((E^\bullet,D)\) 为光滑流形上的微分复形。其第 \(q\) 次光滑上同调为
	\[
		H_D^q(M)=
		\frac{\{u\in C^\infty(E^q):D_qu=0\}}
		{D_{q-1}C^\infty(E^{q-1})}.
	\]
	两个闭截面等价，恰好表示它们之差为 \(D_{q-1}v\)，其中 \(v\) 光滑。
\end{definition}
\symidx{\(H_D^q(M),\mathcal H_D^q\)：微分复形上同调及调和截面空间}

\begin{theorem}\label{g3:thm:elliptic-complex-harmonic-representatives}
	设 \(M\) 为维数 \(n\ge1\) 的紧致无边界光滑流形，\((E^\bullet,D)\) 为带光滑 Hermitian 度量与正密度的一阶椭圆复形。映射
	\[
		\mathcal H_D^q\longrightarrow H_D^q(M),\qquad h\longmapsto[h],
	\]
	是同构，其逆将任意闭代表 \(u\) 送到 \(H_qu\)。因此 \(H_D^q(M)\) 有限维，每个上同调类有唯一的调和代表；该代表也是这一类中唯一的 \(L^2\) 范数最小者。
\end{theorem}
\begin{proof}
	调和截面闭合，因此映射有定义。若 \(h=D_{q-1}v\) 又调和，则
	\(\|h\|_2^2=(h,D_{q-1}v)=(D_{q-1}^*h,v)=0\)，故映射单射。若 \(D_qu=0\)，Green 交换关系给出 \(D_qG_qu=G_{q+1}D_qu=0\)，于是
	\[
		u=H_qu+D_{q-1}D_{q-1}^*G_qu.
	\]
	第二项正合，所以 \([u]=[H_qu]\)，映射满射。更换代表增加 \(D_{q-1}v\) 时，调和投影不变，因而逆映射与代表选择无关。有限维性随 \(\mathcal H_D^q\) 的有限维性而来。

	若 \(h\) 为该类的调和代表，任何其他代表都写为 \(h+D_{q-1}v\)，且两项正交。因此
	\[
		\|h+D_{q-1}v\|_2^2=\|h\|_2^2+\|D_{q-1}v\|_2^2.
	\]
	最小值恰在正合分量为零时达到，即代表等于 \(h\)。
\end{proof}

\begin{corollary}\label{g3:cor:canonical-complex-primitive}
	设 \(M\) 为维数 \(n\ge1\) 的紧致无边界光滑流形，\((E^\bullet,D)\) 为带光滑 Hermitian 度量与正密度的一阶椭圆复形。光滑闭截面 \(u\in C^\infty(E^q)\) 正合当且仅当 \(H_qu=0\)。此时下述原像及估计对每个整数 \(k\ge0\) 成立：
	\[
		v=D_{q-1}^*G_qu,\qquad D_{q-1}v=u,
		\qquad \|v\|_{H^{k+1}}\le C_k\|u\|_{H^k}\quad(k\ge0).
	\]
	该 \(v\) 正交于所有光滑的 \(D_{q-1}\)-闭截面，是光滑原像中唯一的 \(L^2\) 范数最小者。
\end{corollary}
\begin{proof}
	正合性判别及原像公式已由前一定理给出。Green 算子将 \(H^k\) 有界送到 \(H^{k+2}\)，再施加一阶 \(D_{q-1}^*\) 即得估计。若 \(D_{q-1}w=0\)，则
	\((v,w)=(G_qu,D_{q-1}w)=0\)。任何其他光滑原像为 \(v+w\)，由正交性，平方范数增加 \(\|w\|_2^2\)，从而得到唯一最小性。
\end{proof}

改变正密度或丛度量会改变形式伴随、调和空间与最小范数代表，却不会改变 \(H_D^q(M)\) 的代数定义。若两种选择对应调和投影 \(H_q^{(0)},H_q^{(1)}\)，则
\[
	\mathcal H_{D,(0)}^q\longrightarrow\mathcal H_{D,(1)}^q,
	\qquad h\longmapsto H_q^{(1)}h
\]
是同构：它保持同一个上同调类，逆映射为另一方向的调和投影。这个识别说明有限维障碍属于复形，而代表的形状记录着度量选择。

\subsection{de Rham 和 Dolbeault 的共同机制}
\label{g3:subsec:1304:05}
先看实微分形式的复化。记 \(\Omega^q(M;\mathbb C)=C^\infty(\Lambda^q\mathrm T^*M\otimes\mathbb C)\)，局部写
\[
	\mathrm d\left(\sum_I f_I\,\mathrm dx^I\right)
	=\sum_{j,I}\partial_jf_I\,\mathrm dx^j\wedge\mathrm dx^I.
\]
链式法则保证不同坐标中的表达式相容，混合偏导的对称性与外乘的反对称性给出 \(\mathrm d^2=0\)。这是 de Rham 复形。它的符号是否正合，只需一条外代数恒等式。

\begin{lemma}\label{g3:lem:wedge-contraction-symbol}
	设 \(V\) 为有限维实内积空间或复 Hermitian 空间，在其外代数上使用诱导内积。对 \(\eta\in V^*\)，令 \(\varepsilon_\eta\) 为左外乘 \(\eta\)，\(\varepsilon_\eta^*\) 为其伴随，则
	\[
		\varepsilon_\eta^*\varepsilon_\eta+
		\varepsilon_\eta\varepsilon_\eta^*=|\eta|^2I.
	\]
	若 \(\eta\ne0\)，外乘 \(\eta\) 的有限外代数序列处处正合。
\end{lemma}
\begin{proof}
	取正交或酉余标架，使 \(\eta=|\eta|e^1\)。每个形式唯一写成 \(e^1\wedge\alpha+\beta\)，其中 \(\alpha,\beta\) 都不含 \(e^1\)。两部分正交；\(\varepsilon_\eta\) 将 \(\beta\) 送到 \(|\eta|e^1\wedge\beta\)，并消去第一部分，其伴随则把 \(e^1\wedge\alpha\) 送到 \(|\eta|\alpha\)，并消去第二部分。逐项作用即得恒等式。若 \(\varepsilon_\eta u=0\)，则
	\(u=|\eta|^{-2}\varepsilon_\eta\varepsilon_\eta^*u\)，所以 \(u\) 属于前一外乘的像，正合性得证。
\end{proof}

对实余切向量 \(\xi\)，由 \(\partial_j=\mathrm i\mathrm D_j\) 得到
\[
	\sigma_1(\mathrm d)(x,\xi)=\mathrm i\varepsilon_\xi,
	\qquad
	\sigma_1(\mathrm d^*)(x,\xi)=-\mathrm i\varepsilon_\xi^*,
	\qquad
	\sigma_2(\Delta_{\mathrm d})(x,\xi)=|\xi|_g^2I.
\]
因此 de Rham 复形椭圆，关联的 \(\Delta_{\mathrm d}=\mathrm d\mathrm d^*+\mathrm d^*\mathrm d\) 具有非负 Laplace 型主符号。在任意紧致无边界黎曼流形上，光滑 de Rham 上同调已经具有上述唯一调和代表。第十五章将以形式为中心恢复 Hodge 星、配对与曲率公式，进一步解释这些代表的几何性质。

\begin{example}\label{g3:exm:flat-torus-harmonic-forms}
	在平坦环面 \(\mathbb T^n=(\mathbb R/2\pi\mathbb Z)^n\) 上，余标架 \(\mathrm dx^j\) 平行且正交。分部积分得
	\[
		\mathrm d=\sum_j\varepsilon_{\mathrm dx^j}\partial_j,
		\mathrm d^*=-\sum_j\varepsilon_{\mathrm dx^j}^*\partial_j,
		\Delta_{\mathrm d}=-\sum_j\partial_j^2
	\]
	逐分量作用。最后一个等式中，同指标用上一引理，不同指标的外乘与收缩交叉项相消；也可在 \(\mathrm dx^I\) 基底上逐项计算。对 \(u=\sum_Iu_I\mathrm dx^I\) 与自身积分配对，得到
	\[
		(\Delta_{\mathrm d}u,u)=\sum_{I,j}\int_{\mathbb T^n}|\partial_ju_I|^2.
	\]
	故调和形式恰为系数恒定的形式。因此 \(\dim H_{\mathrm d}^q(\mathbb T^n)=\binom nq\)，调和投影就是对每个系数取平均。这是由微分方程直接算出的上同调维数。
\end{example}

复几何中相应的微分只增加反全纯次数。为明确这里需要的条件，设 \(X\) 是复维数 \(n\) 的复流形，即具有双全纯坐标变换的复坐标图册；记 \(\Omega^{p,q}(X)\) 为局部由 \(\mathrm dz^I\wedge\mathrm d\bar z^J\) 张成的光滑形式，\(|I|=p,|J|=q\)。写 \(z^j=x^j+\mathrm iy^j\)，定义
\[
	\bar\partial u
	=\sum_{j,I,J}\frac{\partial u_{I,J}}{\partial\bar z^j}
	\,\mathrm d\bar z^j\wedge\mathrm dz^I\wedge\mathrm d\bar z^J,
	\qquad
	\frac{\partial}{\partial\bar z^j}
	=\frac12\left(\partial_{x^j}+\mathrm i\partial_{y^j}\right).
\]
坐标变换的全纯性使 \(\mathrm dz^I\) 的变换系数被 \(\bar\partial\) 消去，链式法则因而给出全局算子；再次求导时，反全纯混合偏导交换而外乘反交换，故 \(\bar\partial^2=0\)。固定 \(p\) 后得到\term[Dolbeault fuxing]{Dolbeault 复形}[Dolbeault complex]。\footnote{Pierre Dolbeault（1924--2015），法国数学家，研究多复变函数、解析几何与微分形式上同调；\(\bar\partial\) 上同调及其与全纯对象的联系以其命名。}

\begin{proposition}\label{g3:prop:dolbeault-elliptic-symbol}
	设 \(X\) 为复维数 \(n\ge1\) 的复流形，给定 Hermitian 度量，其相应实黎曼度量记为 \(g\)。固定 \(0\le p\le n\)，则
	\[
		0\longrightarrow\Omega^{p,0}(X)\xrightarrow{\bar\partial}
		\Omega^{p,1}(X)\longrightarrow\cdots
		\xrightarrow{\bar\partial}\Omega^{p,n}(X)\longrightarrow0
	\]
	为椭圆复形。采用由 \(g\) 复化诱导的形式内积，对实 \(\xi\in\mathrm T_x^*X\)，有
	\[
		\sigma_1(\bar\partial)(x,\xi)=\mathrm i\varepsilon_{\xi^{0,1}},
		\qquad
		\sigma_2(\Delta_{\bar\partial})(x,\xi)
		=|\xi^{0,1}|^2I=\frac12|\xi|_g^2I,
	\]
	其中 \(\Delta_{\bar\partial}=\bar\partial\bar\partial^*+\bar\partial^*\bar\partial\)。若 \(X\) 紧且无边界，则
	\[
		H_{\bar\partial}^{p,q}(X)
		=\frac{\ker(\bar\partial:\Omega^{p,q}\to\Omega^{p,q+1})}
		{\bar\partial\Omega^{p,q-1}}
	\]
	有限维，并与 \(\ker\Delta_{\bar\partial}\) 自然同构。
\end{proposition}
\begin{proof}
	在坐标中，实协向量 \(\xi=\sum_j(a_j\,\mathrm dx^j+b_j\,\mathrm dy^j)\) 的反全纯部分为
	\[
		\xi^{0,1}=\frac12\sum_j(a_j+\mathrm i b_j)\,\mathrm d\bar z^j.
	\]
	前述 \(\bar\partial\) 表达式立刻给出其主符号。实 \(\xi\) 的两种类型分量互为共轭；Hermitian 度量下二者正交且等长，所以 \(|\xi^{0,1}|^2=|\xi|_g^2/2\)，特别地非零实 \(\xi\) 的反全纯部分非零。把\cref{g3:lem:wedge-contraction-symbol}用于反全纯外代数，固定的全纯外代数只作为有限维系数空间，便得到符号正合性。若把全纯因子统一放在前面，外乘可能带有固定的 \((-1)^p\)，它不改变正合性，也在与伴随相乘时抵消。

	由伴随符号公式，\(\sigma_1(\bar\partial^*)=-\mathrm i\varepsilon_{\xi^{0,1}}^*\)。外乘收缩恒等式于是给出所述 Laplace 符号。最后对紧致无边界的 \(X\) 应用\cref{g3:thm:elliptic-complex-harmonic-representatives}，即得有限维性与唯一调和代表。
\end{proof}

这里复坐标的存在保证 \(\bar\partial^2=0\)，Hermitian 度量则用于定义伴随和代表元。所得调和分解适用于一般紧复流形；将 \(\mathrm d\)-调和空间进一步按复类型分解，还需要另外的几何恒等式。这样，de Rham 与 Dolbeault 的共同部分明确落在“符号正合、非负椭圆 Laplace、Green 分解”这三步中。
% 纲要条目排为编号小节标题，并进入目录与书签；说明段仍保留于纲要注释中。

% !TeX root = ../../Book3.tex
% 纲要标题：### 13.5 Banach 微分与非线性隐函数方法
% 纲要位置：几何学计划纲要.md，第 2642 行。
% 正文状态：初稿；尚未正式审核。

% 纲要内容开始
% 
% 本节是非线性几何的共同工具，独立于第24章的 Yamabe 专题；逆函数、分裂满射和 Fredholm 约化三种情形分别列明假设。
% 
% * Fréchet 微分、连续可微性与链式法则
% * Banach 空间逆函数与隐函数定理的定量证明
% * 分裂满射的零点集及有限维核的局部参数
% * Fredholm 线性化的核、像和余核分解
% * Lyapunov–Schmidt 有限维约化与障碍映射
% * 压缩映射、二次余项和参数一致性
% 
% 纲要内容结束

\section{Banach 微分与非线性隐函数方法}
\label{g3:sec:1305}

线性椭圆理论把核与余核集中到有限维空间。面对非线性方程，还要回答两个问题：线性化的逆怎样控制真实解，以及无法反演的有限维方向留下什么方程。本节先在一般 Banach 空间中建立微分和定量求解，再将满射与 Fredholm 两种线性化分别转化为局部参数和障碍映射。后续几何应用只需核对具体算子的函数空间、微分与估计。

\subsection{Fréchet 微分、连续可微性与链式法则}
\label{g3:subsec:1305:01}
以下以实 Banach 空间为基本对象；复空间上的一般几何非线性映射也先按实空间微分。记 \(\mathcal L(X,Y)\) 为有界线性算子空间，赋予算子范数。它是 Banach 空间：算子范数 Cauchy 列在每个向量上收敛，逐点极限线性且有界，再取单位球上确界便得算子范数收敛。

\begin{definition}\label{g3:def:frechet-derivative}
	设 \(X,Y\) 为 Banach 空间，\(U\subset X\) 开、\(F:U\to Y\)，并取 \(x\in U\)。若存在 \(A\in\mathcal L(X,Y)\)，使
	\[
		F(x+h)=F(x)+Ah+R_x(h),\qquad
		\frac{\|R_x(h)\|_Y}{\|h\|_X}\longrightarrow0,
	\]
	则称 \(F\) 在 \(x\) 处 Fréchet 可微，记 \(\mathrm DF(x)=A\)。若它处处可微，且 \(x\mapsto\mathrm DF(x)\) 在算子范数中连续，则称 \(F\in C^1(U,Y)\)。递归微分定义 \(C^r\)，其 \(r\) 阶微分取值于连续 \(r\) 重线性映射空间。
\end{definition}
\term[Frechet weifen]{Fréchet 微分}[Fréchet derivative]要求余项对所有趋零方向统一变小。\footnote{René Maurice Fréchet（1878--1973），法国数学家，研究抽象空间、点集拓扑与泛函分析。}导数的唯一性来自固定 \(v\)、令 \(h=tv\)：两个候选导数之差作用于 \(v\) 必为零。只在每个固定方向分别求导，尚未给出上述统一余项；这是把形式线性化用于函数空间时必须检查的条件。

\begin{proposition}\label{g3:prop:banach-calculus}
	设 \(X,Y,Z\) 为 Banach 空间，\(F\) 与 \(G\) 在相应开集上可微，且复合有定义，则 Fréchet 微分满足链式法则
	\[
		\mathrm D(G\circ F)(x)=\mathrm DG(F(x))\,\mathrm DF(x).
	\]
	若 \(F\in C^1(U,Y)\)，线段 \([x,x+h]\subset U\)，则
	\begin{equation}\label{g3:eq:banach-fundamental-calculus}
		F(x+h)-F(x)=\int_0^1\mathrm DF(x+th)h\,\mathrm dt.
	\end{equation}
	特别地，在凸集上 \(\sup\|\mathrm DF\|\le L\) 蕴含 \(F\) 为 \(L\)-Lipschitz。连续双线性映射满足乘积法则。
\end{proposition}
\begin{proof}
	把 \(F(x+h)-F(x)=Ah+o(\|h\|)\) 代入 \(G\) 在 \(F(x)\) 处的展开。该增量为 \(O(\|h\|)\)，故第二个余项也是 \(o(\|h\|)\)，得到链式法则。

	连续的 Banach 值函数在紧区间上一致连续，其 Riemann 和为 Cauchy 列；完备性定义积分，并给出 \(\|\int v\|\le\int\|v\|\)。对任意 \(\ell\in Y^*\)，一元微积分应用于 \(\ell(F(x+th))\)，得到式\eqref{g3:eq:banach-fundamental-calculus}两边经 \(\ell\) 后相等。第十一章的 Hahn--Banach 分离结论于是给出向量恒等式。取范数便得 Lipschitz 界。双线性映射 \(B\) 的展开中，额外项 \(B(h,k)\) 为 \(O(\|h\|\|k\|)\)，从而得到乘积法则。
\end{proof}

例如在一个 Banach 代数中，\(F(u)=u^2\) 的导数是 \(\mathrm DF(u)h=uh+hu\)，余项恰为 \(h^2\)。代数不交换时不能把它改写成 \(2uh\)。这说明线性化既要保留乘法次序，也要在实际目标范数中估计余项。

\subsection{Banach 空间逆函数与隐函数定理的定量证明}
\label{g3:subsec:1305:02}
\begin{lemma}[压缩映射原理]\label{g3:lem:banach-contraction}
	设 \(K\) 为非空完备度量空间，\(T:K\to K\) 满足
	\(d(Tx,Ty)\le qd(x,y)\)，其中 \(0\le q<1\)。则 \(T\) 有唯一不动点 \(x_*\)，且从任意 \(x_0\in K\) 出发的迭代 \(x_{j+1}=T(x_j)\) 满足
	\[
		d(x_j,x_*)\le\frac{q^j}{1-q}d(x_1,x_0).
	\]
\end{lemma}
\begin{proof}
	逐次应用压缩性，得 \(d(x_{j+1},x_j)\le q^j d(x_1,x_0)\)。对相邻距离求和，几何级数使 \((x_j)\) 为 Cauchy 列，完备性给出极限 \(x_*\)。\(T\) 连续，所以 \(Tx_*=x_*\)；尾和给出误差界。若 \(y_*\) 也为不动点，则 \(d(x_*,y_*)\le qd(x_*,y_*)\)，故二者相同。
\end{proof}

\begin{theorem}[定量逆函数定理]\label{g3:thm:quantitative-banach-inverse}
	设 \(X,Y\) 为 Banach 空间，\(F\in C^1(U,Y)\)，\(\overline B_r(x_0)\subset U\)，且 \(A=\mathrm DF(x_0):X\to Y\) 有有界逆。若
	\begin{equation}\label{g3:eq:inverse-small-derivative}
		\sup_{x\in\overline B_r(x_0)}
		\|I-A^{-1}\mathrm DF(x)\|\le q<1,
	\end{equation}
	则 \(F\) 在该球上单射。每个满足
	\(\|A^{-1}(y-F(x_0))\|<(1-q)r\) 的 \(y\)，在开球中有唯一原像 \(x=G(y)\)，并且
	\[
		\|G(y)-G(z)\|\le\frac{\|A^{-1}\|}{1-q}\|y-z\|,
		\qquad \mathrm DG(y)=\mathrm DF(G(y))^{-1}.
	\]
	\(F\) 在 \(x_0\) 附近是 \(C^1\) 微分同胚；若 \(F\in C^r\)，\(1\le r\le\infty\)，则局部逆也属于 \(C^r\)。
\end{theorem}
\begin{proof}
	令 \(T_y(x)=x-A^{-1}(F(x)-y)\)。由上一命题和式\eqref{g3:eq:inverse-small-derivative}，它在闭球中为 \(q\)-压缩，且
	\[
		\|T_y(x)-x_0\|
		\le q\|x-x_0\|+\|A^{-1}(y-F(x_0))\|\le r.
	\]
	压缩原理给出唯一不动点，正是 \(F(x)=y\) 的解。严格的右端条件及同一估计使 \(\|x-x_0\|<r\)。对两个解相减，得
	\[
		(1-q)\|x-x'\|\le\|A^{-1}(F(x)-F(x'))\|;
	\]
	这同时证明单射与逆的 Lipschitz 界。

	算子 \(A^{-1}\mathrm DF(x)=I-B_x\) 满足 \(\|B_x\|\le q\)。几何级数 \(\sum_{j\ge0}B_x^j\) 在算子范数中收敛，与 \(I-B_x\) 相乘后为恒等，所以每个 \(\mathrm DF(x)\) 可逆，逆范数至多为 \(\|A^{-1}\|/(1-q)\)。在每个内点重新取小闭球，重复不动点构造，可知 \(F(B_r(x_0))\) 为开集。

	写 \(x=G(y)\)、\(h=G(y+k)-G(y)\)。可微性给出
	\(k=\mathrm DF(x)h+o(\|h\|)\)。已得 \(\|h\|=O(\|k\|)\)，故
	\(h=\mathrm DF(x)^{-1}k+o(\|k\|)\)，证明逆的导数公式。逆算子的连续性由
	\[
		B^{-1}-C^{-1}=B^{-1}(C-B)C^{-1}
	\]
	得到；因而 \(G\in C^1\)。进一步，在可逆算子组成的开集上，上式给出
	\(\mathrm D(B\mapsto B^{-1})[H]=-B^{-1}HB^{-1}\)，余项为 \(O(\|H\|^2)\)。反复使用乘积法则，求逆映射光滑。由 \(\mathrm DG=(\mathrm DF\circ G)^{-1}\) 归纳即得高阶结论。
\end{proof}

连续性保证在 \(\mathrm DF(x_0)\) 有界可逆时，缩小 \(r\) 总能满足式\eqref{g3:eq:inverse-small-derivative}。条件中的逆范数与导数变化率决定可用邻域的大小，这比单说“在附近可解”保留了更多可用于参数族的信息。

\begin{corollary}[带参数的隐函数定理]\label{g3:cor:banach-implicit-parameters}
	设 \(Z,X,Y\) 为 Banach 空间，\(F\) 在 \((z_0,x_0)\) 附近为 \(C^r\) 映射，\(r\ge1\)，且 \(\mathrm D_xF(z_0,x_0):X\to Y\) 有有界逆。则在 \((z_0,F(z_0,x_0))\) 附近，方程 \(F(z,x)=y\) 有唯一的近邻解 \(x=\varphi(z,y)\)，且 \(\varphi\in C^r\)。
\end{corollary}
\begin{proof}
	对 \(\mathcal F(z,x)=(z,F(z,x))\) 应用逆函数定理。其导数为块三角矩阵，对 \((\zeta,\eta)\) 的逆作用为
	\((\zeta,(\mathrm D_xF)^{-1}(\eta-\mathrm D_zF\,\zeta))\)，确有有界逆。逆映射的第二分量就是 \(\varphi\)。微分恒等式还给出
	\[
		\mathrm D_z\varphi=-(\mathrm D_xF)^{-1}\mathrm D_zF,\qquad
		\mathrm D_y\varphi=(\mathrm D_xF)^{-1},
	\]
	其中右边在 \((z,\varphi(z,y))\) 处取值。
\end{proof}

\subsection{分裂满射的零点集及有限维核的局部参数}
\label{g3:subsec:1305:03}
\begin{definition}\label{g3:def:split-surjection}
	设 \(X,Y\) 为 Banach 空间。有界线性算子 \(L:X\to Y\) 称为\term[fenlie manshe]{分裂满射}[split surjection]，若存在有界线性 \(Q:Y\to X\) 使 \(LQ=I_Y\)。这里 \(Q\) 是一个有界右逆。
\end{definition}
由 \((QL)^2=QL\)，有界投影 \(QL\) 给出
\[
	X=\ker L\oplus X_1,\qquad X_1=\operatorname{im}Q.
\]
两部分均闭，投影分别为 \(I-QL\) 与 \(QL\)，而 \(L|_{X_1}\) 的逆为 \(Q\)。这个分解把保持方程不变的一阶方向与实际被方程控制的方向分开。

\begin{theorem}\label{g3:thm:split-regular-level}
	设 \(X,Y\) 为 Banach 空间，\(U\subset X\) 开，\(F:U\to Y\) 为 \(C^r\)，\(r\ge1\)，\(F(x_0)=y_0\)，且 \(L=\mathrm DF(x_0)\) 为分裂满射。按上式取 \(K=\ker L\)、\(X_1=\operatorname{im}Q\)。存在零点邻域上的 \(C^r\) 映射 \(\phi:K\to X_1\)，使 \(x_0\) 附近的水平集恰为
	\[
		F^{-1}(y_0)=\{x_0+k+\phi(k):k\in K\text{ 充分小}\},
		\qquad\phi(0)=0,\quad\mathrm D\phi(0)=0.
	\]
	因此该水平集是以 \(K\) 为模型的局部 Banach 子流形；若 \(\dim K=N<\infty\)，它在此处有 \(N\) 个局部参数。
\end{theorem}
\begin{proof}
	将 \(F(x_0+k+z)=y_0\) 看成对 \(z\in X_1\) 的方程。其 \(z\) 导数为有界同构 \(L|_{X_1}\)，带参数隐函数定理给出唯一 \(z=\phi(k)\)。每个近邻 \(x\) 都有唯一的 \(K\oplus X_1\) 分解，因此没有遗漏其他解。对等式在零点微分，得
	\(L(k+\mathrm D\phi(0)k)=0\)。\(Lk=0\)，且 \(L|_{X_1}\) 单射，所以 \(\mathrm D\phi(0)=0\)。图的微分为 \(h\mapsto h+\mathrm D\phi(k)h\)，并把模型空间识别为该点线性化的核。
\end{proof}

在 Hilbert 空间中可以借助正交补寻找右逆；在一般 Banach 空间中则必须保留“有界补空间”条件。下一小节将证明，有有限维核且闭值域的 Fredholm 算子总有本节所需的分解。

\begin{example}\label{g3:exm:banach-sphere-regular-level}
	在实 Hilbert 空间 \(H\) 中，令 \(F(u)=\|u\|^2\)。当 \(\|u_0\|=1\) 时，\(\mathrm DF(u_0)h=2\langle h,u_0\rangle\)，右逆为 \(Q(t)=tu_0/2\)，核是 \(u_0^\perp\)。单位球面在 \(u_0\) 附近由
	\[
		k\longmapsto k+\sqrt{1-\|k\|^2}\,u_0,
		\qquad k\in u_0^\perp,\quad\|k\|<1,
	\]
	参数化。有限维的球面图形在这里直接变成 Hilbert 空间中的水平集模型。
\end{example}

\subsection{Fredholm 线性化的核、像和余核分解}
\label{g3:subsec:1305:04}
要从一个连续线性双射取得有界逆，需要使用完备性。下面把这一功能明确补齐。

\begin{lemma}[开映射定理]\label{g3:lem:banach-open-mapping}
	Banach 空间之间的有界线性满射 \(L:X\to Y\) 是开映射。特别地，有界线性双射的逆也有界。
\end{lemma}
\begin{proof}
	先证完备度量空间不能是可数个闭的无内点集合 \(F_j\) 之并。这是所需的 Baire 性质。\footnote{René-Louis Baire（1874--1932），法国数学家，研究函数分类、连续性与极限理论。}若这种覆盖成立，可依次选择非空闭球 \(\overline B_j\)，使它位于前一球的内部、避开 \(F_j\)，且半径小于 \(2^{-j}\)。选择可能，因为 \(F_j\) 闭且无内点。球心列为 Cauchy 列，其极限属于全部闭球，故不在任何 \(F_j\)，矛盾。

	由满射性，\(Y=\bigcup_{j\ge1}\overline{L(jB_X)}\)，其中 \(B_X\) 为开单位球。Baire 性质给出某个闭集包含一个开球。将该球内的两个点相减，利用线性与闭包的相容性，再缩放，得到某个 \(c>0\) 满足
	\[
		B_Y(0,c)\subset\overline{L(B_X)}.
	\]
	现在去掉闭包。若 \(\|y\|<c/2\)，先选 \(\|x_1\|<1/2\)，使 \(\|y-Lx_1\|<c/4\)；再选 \(\|x_2\|<1/4\)，使新余项小于 \(c/8\)，如此继续。于是 \(\sum x_j\) 在 \(X\) 中收敛，和 \(x\) 满足 \(\|x\|<1\)、\(Lx=y\)。因此 \(B_Y(0,c/2)\subset L(B_X)\)。平移与缩放证明开映射；双射时，同一包含给出 \(\|L^{-1}y\|\le(2/c)\|y\|\)。
\end{proof}

\begin{proposition}\label{g3:prop:banach-fredholm-splitting}
	设 \(X,Y\) 为 Banach 空间，有界线性 \(L:X\to Y\) 为 Fredholm 算子，即 \(K=\ker L\) 有限维、\(R=\operatorname{im}L\) 闭且 \(Y/R\) 有限维。则存在有界投影及闭分解
	\[
		X=K\oplus X_1,\qquad Y=R\oplus C,
	\]
	其中 \(C\) 有限维，\(L_1=L|_{X_1}:X_1\to R\) 有有界逆 \(Q\)。若 \(R=Y\)，则 \(L\) 为分裂满射。
\end{proposition}
\begin{proof}
	在 \(K\) 中选基 \(e_1,\ldots,e_a\)，其坐标泛函连续。由 Hahn--Banach 将它们延拓为 \(\ell_i\in X^*\)，则 \(\Pi_Kx=\sum_i\ell_i(x)e_i\) 是有界投影，取 \(X_1=\ker\Pi_K\)。因为 \(R\) 闭，商空间 \(Y/R\) 的商范数是范数；它有限维，所以任何线性坐标与截面都连续。选择一组商基的代表，张成 \(C\)，便得到连续截面 \(S:Y/R\to Y\) 及投影 \(\Pi_C=S\circ\pi\)。\(I-\Pi_C\) 投影到 \(R\)。\(L_1\) 为闭子空间间的有界双射，开映射定理使其逆有界。
\end{proof}

椭圆情形中的 \(K\) 和余核由光滑解表示，可以使用前面构造的正交投影；上述论证说明一般 Banach 线性化也有足够的分解。不同补空间会改变局部坐标与障碍映射的具体公式，方程的真实零点集合则由这些坐标互相对应。

\subsection{Lyapunov–Schmidt 有限维约化与障碍映射}
\label{g3:subsec:1305:05}
\term[Lyapunov Schmidt yuehua]{Lyapunov--Schmidt 约化}[Lyapunov--Schmidt reduction]先解值域中的方程，再把剩余条件放到余核。\footnote{Aleksandr Mikhailovich Lyapunov（1857--1918），俄国数学家，研究动力系统稳定性、数学物理与概率论。Erhard Schmidt（1876--1959），德国数学家，研究积分方程与 Hilbert 空间理论。}

\begin{theorem}\label{g3:thm:lyapunov-schmidt-reduction}
	设 \(X,Y\) 为 Banach 空间，\(U\subset X\) 开，\(F:U\to Y\) 为 \(C^r\)，\(r\ge1\)，\(F(0)=0\)，且 \(L=\mathrm DF(0)\) 为 Fredholm 算子。取上一命题的分解，记 \(\Pi_R=I-\Pi_C\)。存在 \(C^r\) 映射 \(\psi:K\supset V\to X_1\)，满足
	\[
		\psi(0)=0,\quad\mathrm D\psi(0)=0,\qquad
		\Pi_RF(k+\psi(k))=0.
	\]
	定义\term[zhangai yingshe]{障碍映射}[obstruction map]
	\[
		\kappa(k)=\Pi_CF(k+\psi(k))\in C.
	\]
	则零点附近，\(F(x)=0\) 与有限维方程 \(\kappa(k)=0\) 等价，具体对应为 \(x=k+\psi(k)\)。并且 \(\kappa(0)=0\)、\(\mathrm D\kappa(0)=0\)。
\end{theorem}
\begin{proof}
	考虑 \(\Pi_RF(k+z)=0\)，未知量为 \(z\in X_1\)，参数为 \(k\in K\)。在原点对 \(z\) 的微分为 \(L_1\)，是有界同构，故隐函数定理给出唯一的 \(z=\psi(k)\)。对 \(k\) 微分并用 \(Lk=0\)，得到 \(\mathrm D\psi(0)=0\)。原方程等价于 \(R\) 与 \(C\) 两个投影同时为零；前者已经解出，后者正是 \(\kappa(k)=0\)。最后 \(\Pi_CL=0\)，所以 \(\mathrm D\kappa(0)=0\)。
\end{proof}

若 \(F\in C^2\)，则 \(\psi(k)=O(\|k\|^2)\)，而首个可能非零的障碍为
\begin{equation}\label{g3:eq:quadratic-obstruction}
	\kappa(k)=\frac12\Pi_C\mathrm D^2F(0)[k,k]+o(\|k\|^2).
\end{equation}
这是因为 \(\Pi_CL\psi=0\)，且把 \(k+O(\|k\|^2)\) 代入二阶项只改变更高阶余项。若 \(C=0\)，约化没有剩余方程，恢复正则水平集定理；若 \(C\ne0\)，即使线性化有非零核，也未必有经过原点的非平凡解族。

\begin{example}\label{g3:exm:quadratic-obstruction-isolated}
	设 \(H\) 为实 Hilbert 空间，\(F:\mathbb R\oplus H\to\mathbb R\oplus H\) 为 \(F(t,z)=(t^2,z)\)。在原点，核与余核各一维、指数为零，值域方程给出 \(z=0\)，障碍是 \(\kappa(t)=t^2\)。因此原点仍是孤立零点。改为 \(F_\lambda(t,z)=(t^2-\lambda,z)\)，则 \(\lambda<0\) 无实解，\(\lambda=0\) 有一个，\(\lambda>0\) 有两个近邻解。有限维障碍保留了线性化没有记录的分支信息。
\end{example}

参数化方程可把外参数与 \(k\) 一起放进隐函数定理，得到 \(\psi(k,a)\)、\(\kappa(k,a)\)。这里核、值域与余核首先取自固定基点的线性化；近邻方程的变化进入非线性项，而非在核维数可能跳跃时任意选择一族核基。

\subsection{压缩映射、二次余项和参数一致性}
\label{g3:subsec:1305:06}
\begin{proposition}\label{g3:prop:uniform-quadratic-solve}
	设 \(X,Y\) 为 Banach 空间，\(F:X\supset\overline B_r(0)\to Y\) 在该球邻域为 \(C^1\)，写
	\[
		F(x)=b+Ax+N(x),\qquad N(0)=0,
	\]
	其中 \(A\) 有有界逆，\(\|A^{-1}\|\le M\)，且
	\[
		\|\mathrm DN(x)\|\le L\|x\|\quad(\|x\|\le r).
	\]
	若 \(MLr\le1/2\)、\(M\|b\|\le r/2\)，则 \(F(x)=0\) 在闭球中有唯一解，并有
	\[
		\|x\|\le2M\|b\|,\qquad
		\|x+A^{-1}b\|\le\frac{ML}{2}\|x\|^2.
	\]
	所有满足同一组 \(M,L,r\) 的方程共享这些结论。
\end{proposition}
\begin{proof}
	令 \(T(x)=-A^{-1}b-A^{-1}N(x)\)。微积分公式给出
	\[
		\|N(x)\|\le\frac L2\|x\|^2,\qquad
		\|N(x)-N(y)\|\le Lr\|x-y\|
	\]
	于闭球内成立。因此 \(T\) 的压缩常数至多为 \(1/2\)，且
	\(\|T(x)\|\le r/2+MLr^2/2\le3r/4\)。压缩原理给出唯一解。由
	\(\|x\|\le M\|b\|+MLr\|x\|/2\le M\|b\|+\|x\|/2\)
	得到第一界；固定点恒等式及 \(N\) 的二次界给出第二界。
\end{proof}

若 \(F\in C^2\) 且 \(\|\mathrm D^2F\|\le L\)，取 \(A=\mathrm DF(0)\) 后便有上述导数界。一般 \(C^1\) 映射的余项只是 \(o(\|x\|)\)；二次余项需要导数具有局部 Lipschitz 控制，不能由一阶连续可微性自动获得。

\ProseParagraphBreak

参数一致性还给出解的连续依赖。若 \(T_a\) 都将同一闭球映入自身、压缩常数统一为 \(q<1\)，其不动点满足
\[
	\|x_a-x_b\|\le\frac1{1-q}
	\|T_a(x_b)-T_b(x_b)\|.
\]
因此右侧关于参数的连续性、Lipschitz 性分别传给解。若方程对 Banach 参数为 \(C^r\)，则由带参数隐函数定理得到 \(C^r\) 依赖。实际应用中须同时控制线性逆的范数与非线性导数的变化率；当逆范数因小特征值趋于无穷时，可解邻域也可能缩小。
% 纲要条目排为编号小节标题，并进入目录与书签；说明段仍保留于纲要注释中。

% !TeX root = ../../Book3.tex
% 纲要标题：### 13.6 Sobolev 映射空间与非线性算子
% 纲要位置：几何学计划纲要.md，第 2653 行。
% 正文状态：初稿；尚未正式审核。

% 纲要内容开始
% 
% 函数空间、微分阶数、系数正则性及参数范围同时固定；不对任意 Sobolev 指数无条件宣称复合算子光滑。
% 
% * Sobolev 连续嵌入与复合定理
% * 指数图册、丛值截面及局部坐标转换
% * Sobolev 乘法、复合及微分损失的估计
% * 非线性微分算子的可微性和目标空间选择
% 
%   曲面上的 Cauchy–Riemann 模块固定 $p>2$，使用 $W^{1,p}\to L^p$ 的映射；高阶族与规范群作用另行核验所需指数，不将微分损失隐藏在“光滑作用”中。
% 
% * 椭圆自举、先验界与弱解的正则性恢复
% * 直接方法、约束临界点与规范对称的区别
% 
% ---
% 
% 纲要内容结束

\section{Sobolev 映射空间与非线性算子}
\label{g3:sec:1306}

要把上一节用于几何方程，必须先选定一个能容纳未知对象的 Banach 空间，并证明非线性表达式在该空间中可微。连续嵌入控制映射的像，乘法控制导数展开，目标空间则记录方程损失的导数。本节按这个顺序建立 Sobolev 映射图册与所需非线性算子，并用半线性方程和曲面上的 Cauchy--Riemann 方程说明正则性怎样恢复。

\subsection{Sobolev 连续嵌入与复合定理}
\label{g3:subsec:1306:01}
以下固定紧、光滑、无边界源流形 \(M^m\)，\(m\ge1\)，整数 \(k\ge1\)、\(1<p<\infty\)，并假定 \(kp>m\)。第十一章给出
\[
	W^{k,p}(M;E)\longrightarrow C^0(M;E),\qquad
	\|uv\|_{W^{k,p}}\le C\|u\|_{W^{k,p}}\|v\|_{W^{k,p}},
\]
其中 \(E\) 为固定光滑有限秩丛，第二式用于相容的逐点乘法。选任意 \(0<\alpha<\min\{1,k-m/p\}\)，还可得到 \(C^{0,\alpha}\) 控制。因而充分小的 Sobolev 邻域同时是逐点小的邻域。

\begin{lemma}\label{g3:lem:sobolev-multiple-products}
	设 \(M^m\) 为紧光滑无边界流形，\(m\ge1\)、\(k\ge1\) 为整数，\(1<p<\infty\)、\(kp>m\)。若正整数 \(j_1,\ldots,j_a\) 的和不超过 \(k\)，则在固定坐标片及截断下，任意相应导数乘积满足
	\[
		\left\|\prod_{i=1}^a\mathrm D^{j_i}u_i\right\|_{L^p}
		\le C\prod_{i=1}^a\|u_i\|_{W^{k,p}}.
	\]
	取向量或张量的有限次收缩也有同样估计。
\end{lemma}
\begin{proof}
	\(a=1\) 直接成立。\(a\ge2\) 时每个 \(j_i<k\)。令
	\(b_i=\max\{1/p-(k-j_i)/m,0\}\)。若有 \(r\ge2\) 个正的 \(b_i\)，则
	\[
		\sum b_i\le\frac rp-\frac{rk-k}{m}
		=\frac1p+(r-1)\left(\frac1p-\frac km\right)<\frac1p.
	\]
	只有一个正数时也严格小于 \(1/p\)，全为零时同样成立。因此可选有限 \(q_i\ge p\)，使 \(1/q_i>b_i\) 且 \(\sum_i1/q_i=1/p\)。第十一章的高阶嵌入给出 \(\mathrm D^{j_i}u_i\in L^{q_i}\)，包括临界时任取有限指数的情形。Hölder 不等式给出结论；有限分量相加处理收缩。
\end{proof}

\begin{theorem}\label{g3:thm:sobolev-composition-smooth}
	设 \(F:M\times O\to\mathbb R^b\) 光滑，\(O\subset\mathbb R^a\) 开，且 \(kp>m\)。则
	\[
		\mathcal U=\{u\in W^{k,p}(M;\mathbb R^a):u(M)\subset O\}
	\]
	为开集，逐点复合 \(\mathcal N_F(u)(x)=F(x,u(x))\) 定义光滑映射
	\(\mathcal U\to W^{k,p}(M;\mathbb R^b)\)。其 \(r\) 阶微分为
	\begin{equation}\label{g3:eq:sobolev-composition-derivative}
		\mathrm D^r\mathcal N_F(u)[h_1,\ldots,h_r](x)
		=\partial_y^rF(x,u(x))[h_1(x),\ldots,h_r(x)].
	\end{equation}
	在像位于固定紧集邻域、Sobolev 范数有界的充分小邻域内，各阶所需算子范数和余项常数只依赖相应有限阶导数界。
\end{theorem}
\begin{proof}
	\(u(M)\) 紧且在 \(O\) 内，与其补集有正距离；连续嵌入证明 \(\mathcal U\) 开。先对光滑 \(u\) 求坐标导数。重复链式法则后，每项是某个有界的 \(F\) 的偏导在 \((x,u(x))\) 取值，乘以
	\(\mathrm D^{j_1}u\cdots\mathrm D^{j_a}u\)，其中 \(j_i\ge1\)、\(\sum j_i\le k\)。上一引理控制这些乘积，零阶项由紧像控制。因此在 \(\|u\|_{W^{k,p}}\le R\) 的集合上，有 \(\|F(\cdot,u)\|_{W^{k,p}}\le C_R\)。

	对两个充分接近的 \(u,v\)，逐项相减。乘积用望远镜展开，每项有一个 \(u-v\) 的导数因子；系数差用 \(F\) 的再高一阶导数及 \(\|u-v\|_\infty\) 控制。得到
	\begin{equation}\label{g3:eq:sobolev-composition-lipschitz}
		\|F(\cdot,u)-F(\cdot,v)\|_{W^{k,p}}
		\le C_R\|u-v\|_{W^{k,p}}.
	\end{equation}
	所说充分接近保证两像之间的直线段仍位于同一个紧邻域；一般开集无需假定凸。光滑稠密性及连续嵌入使光滑逼近的像保持在这个邻域内。式\eqref{g3:eq:sobolev-composition-lipschitz}使复合列在 \(W^{k,p}\) 中收敛，一致收敛又识别其极限为逐点复合，因而上述估计延伸到一般 \(u\)。

	现在证明真正的 Banach 微分。点态微积分给出
	\[
		F(x,u+h)-F(x,u)-\partial_yF(x,u)h
		=\int_0^1\bigl(\partial_yF(x,u+th)-\partial_yF(x,u)\bigr)h\,\mathrm dt.
	\]
	对 \(\partial_yF\) 使用刚证的局部 Lipschitz 估计，再用 \(W^{k,p}\) 代数估计，右端范数至多为 \(C_R\|h\|_{W^{k,p}}^2\)。其积分是连续的 Banach 值积分。故一阶微分确为式\eqref{g3:eq:sobolev-composition-derivative}，并依 \(u\) 连续。把同一论证用于 \(\partial_y^rF\)，结合连续多线性乘法，逐次证明所有高阶微分。每次只使用有限个偏导，最后得到光滑性。
\end{proof}

这一证明适用于有限坐标片上的光滑纤维映射，也适用于光滑边界域上的局部问题；有限覆盖与第十一章的延拓估计给出同样结论。由此可以内在定义\term[Sobolev yingshe kongjian]{Sobolev 映射空间}[Sobolev mapping space] \(W^{k,p}(M,N)\)：它由连续映射 \(u:M\to N\) 组成，在覆盖其图像的有限组源、目标坐标中，各坐标表达属于局部 \(W^{k,p}\)。坐标变换是刚证明的光滑复合，故定义与所选图无关。目标流形 \(N\) 可以非紧；每个固定 \(u\) 的像仍紧，所需常数只在像的一个紧邻域内计算。

\subsection{指数图册、丛值截面及局部坐标转换}
\label{g3:subsec:1306:02}
取目标 \(N\) 上的辅助光滑黎曼度量。若 \(u_0:M\to N\) 光滑，拉回丛 \(u_0^*\mathrm TN\) 是光滑丛。在零截面附近定义
\begin{equation}\label{g3:eq:sobolev-exponential-chart}
	\mathcal E_{u_0}(\xi)(x)=\exp_{u_0(x)}\xi(x),
	\qquad\xi\in W^{k,p}(M;u_0^*\mathrm TN).
\end{equation}
这里用的是逐点指数映射，而不是把一个向量场积分到时间一后得到的自微分同胚。

\begin{lemma}\label{g3:lem:sobolev-map-smooth-density}
	在 \(kp>m\) 时，每个 \(u\in W^{k,p}(M,N)\) 都可由光滑映射逼近，逼近同时在局部 Sobolev 坐标和 \(C^0\) 意义下成立。
\end{lemma}
\begin{proof}
	只须处理紧集 \(K=u(M)\) 的邻域。选择有限个目标坐标 \((V_i,y_i)\) 及支于其中的光滑函数 \(\chi_i\)，使 \(\sum_i\chi_i=1\) 在 \(K\) 的一个邻域成立。各 \(\chi_i y_i\) 在图外补零，定义
	\[
		\iota(q)=\bigl(\chi_i(q),\chi_i(q)y_i(q)\bigr)_i\in\mathbb R^A.
	\]
	在稍小且闭包紧含的邻域上，它是嵌入。确实，若两点的像相同，选在第一点非零的 \(\chi_i\)，则第二点同样非零，且两点的 \(y_i\) 坐标相同。若 \(\mathrm d\iota(v)=0\)，同一坐标中 \(\mathrm d\chi_i(v)=0\) 和 \(\chi_i\mathrm dy_i(v)=0\)，故 \(v=0\)。紧闭包上的连续单射为到其像的同胚，因而取得所需局部嵌入。

	在 \(\iota(K)\) 附近有光滑回缩 \(\pi\) 到所嵌入的子流形。具体地，法丛映射 \((q,\nu)\mapsto\iota(q)+\nu\) 在零截面的导数把切向与法向直和送到 \(\mathbb R^A\)，为同构；有限维逆函数定理给出局部逆，投影到 \(q\) 即得 \(\pi\)。在紧集附近可统一缩小邻域以保证这些逆相容：若任意小的法向量仍产生不同的两个表示，则取趋向紧集的子列，其两个底点距离也趋零，落入同一局部逆的邻域，与唯一性矛盾。

	由上一复合定理，\(\iota\circ u\in W^{k,p}(M;\mathbb R^A)\)。将它在欧氏目标中光滑逼近为 \(v_j\)，同时有 \(v_j\to\iota\circ u\) 一致。因此大 \(j\) 时 \(v_j(M)\) 位于回缩定义域内。光滑映射 \(u_j=\iota^{-1}\circ\pi\circ v_j\) 由复合估计在 \(W^{k,p}\) 及 \(C^0\) 中趋于 \(u\)，证明完成。
\end{proof}

以 Banach 空间的开集为坐标域、过渡映射光滑的 Hausdorff 空间，称为\term[Banach liuxing]{Banach 流形}[Banach manifold]。这里的映射空间用上述 Sobolev 局部坐标拓扑；连续嵌入使不同映射能由不交的 \(C^0\) 邻域分开。

\begin{theorem}\label{g3:thm:sobolev-map-charts}
	设 \(M^m\) 紧、光滑、无边界，\(N\) 为有限维光滑流形，\(k\ge1\)、\(1<p<\infty\)、\(kp>m\)。式\eqref{g3:eq:sobolev-exponential-chart}以光滑参考映射 \(u_0\) 为中心，构成 \(W^{k,p}(M,N)\) 的光滑 Banach 图册。更换目标度量所得图册光滑相容。
\end{theorem}
\begin{proof}
	映射 \((q,v)\mapsto(q,\exp_qv)\) 在零截面的导数可逆。因此在紧集 \(u_0(M)\) 附近，可以统一选择 \(\varepsilon>0\)，使 \(|\xi(x)|<\varepsilon\) 时逐点指数单射且有光滑局部逆 \(\log_{u_0(x)}\)。连续嵌入保证这一条件在截面空间中给出开集。复合定理同时证明 \(\mathcal E_{u_0}\) 及其逆在固定坐标中光滑，所以它确为图。

	若 \(u_0,u_1\) 两图重叠，其过渡映射逐点为
	\[
		\xi(x)\longmapsto
		\log_{u_1(x)}\bigl(\exp_{u_0(x)}\xi(x)\bigr).
	\]
	这是固定光滑丛之间的光滑纤维映射，因而在 \(W^{k,p}\) 中光滑。不同辅助度量之间的过渡具有同样形式。最后，对任意 Sobolev 映射 \(u\)，上一引理给出足够接近的光滑 \(u_0\)。这些逼近的像均在 \(u(M)\) 的同一紧邻域内，指数映射可用半径因此能统一选取。二者足够一致接近，所以 \(u\) 落入 \(u_0\) 的指数图。这一步保证图册覆盖整个映射空间，而不只覆盖光滑映射本身。
\end{proof}

在一个可能不光滑的 \(u\) 处，切向量仍可逐点看成 \(u^*\mathrm TN\) 的 \(W^{k,p}\) 截面；其函数空间结构由附近光滑参考映射的指数微分传递。过渡矩阵属于 \(W^{k,p}\)，乘法估计保证所得范数相容。这样无需把 \(u^*\mathrm TN\) 误称为由 \(u\) 自动给出的光滑拉回丛。

\subsection{Sobolev 乘法、复合及微分损失的估计}
\label{g3:subsec:1306:03}
前面的代数估计控制同阶乘积。微分方程还会出现“高阶系数乘低阶未知量”的组合，需要如下版本。

\begin{proposition}\label{g3:prop:sobolev-multiplier-loss}
	在 \(M^m\) 紧、光滑、无边界，\(m\ge1\)、\(k\ge1\) 为整数、\(kp>m\)、\(1<p<\infty\) 时，对整数 \(0\le l\le k\)，有
	\begin{equation}\label{g3:eq:sobolev-multiplier-loss}
		\|av\|_{W^{l,p}}\le C\|a\|_{W^{k,p}}\|v\|_{W^{l,p}}.
	\end{equation}
	特别地，固定光滑联络给出连续线性微分映射
	\(\nabla^r:W^{k,p}(E)\to W^{k-r,p}((\mathrm T^*M)^{\otimes r}\otimes E)\)，\(0\le r\le k\)。
\end{proposition}
\begin{proof}
	\(l=0\) 由 \(a\in L^\infty\) 直接得到。对 \(l>0\)，先估计最高阶 Leibniz 项 \((\mathrm D^j a)(\mathrm D^{l-j}v)\)。\(j=0\) 仍用 \(L^\infty\cdot L^p\)；若 \(j=l=k\)，则交换两因子使用同一论证。其他 \(1\le j\le l\) 的项，令
	\[
		b=\max\{1/p-(k-j)/m,0\},\qquad
		c=\max\{1/p-j/m,0\}.
	\]
	若二者为正，则 \(b+c=2/p-k/m<1/p\)；若仅一个为正，除已处理的端项外也严格小于 \(1/p\)。可选有限 \(q_1,q_2\) 使 \(1/q_1>b\)、\(1/q_2>c\)、\(1/q_1+1/q_2=1/p\)。高阶 Sobolev 嵌入及 Hölder 不等式控制该乘积。较低总阶的导数同样处理，有限图与逼近完成证明。联络微分的断言直接由局部表达及光滑系数的乘法估计得到。
\end{proof}

若 \(F\) 为固定的光滑纤维映射，式\eqref{g3:eq:sobolev-composition-derivative}表明，零阶复合保持 \(W^{k,p}\) 的阶数；\(r\) 次微分则消耗 \(r\) 阶。于是
\[
	u\longmapsto A(x,u)\nabla^r u
\]
在 \(kp>m\)、\(k\ge r\) 时是到 \(W^{k-r,p}\) 的光滑映射。其微分为
\[
	h\longmapsto A(x,u)\nabla^r h
	+\partial_uA(x,u)[h]\nabla^r u,
\]
由乘子估计，两项落在同一目标空间。

\ProseParagraphBreak

固定目标上的逐点复合与变化源上的重参数化具有不同的导数计数。在圆周上，\(\mathcal R(t,u)(x)=u(x+t)\) 对参数的候选导数是 \(u'\)。若要求它仍取值于 \(W^{k,p}\)，就额外需要 \(u\in W^{k+1,p}\)。因此本节的光滑复合定理用于固定源上的纤维映射；变化坐标、自微分同胚和规范群作用必须按实际公式核验微分损失。

\subsection{非线性微分算子的可微性和目标空间选择}
\label{g3:subsec:1306:04}
记 \(j^ru(x)\) 为截面在 \(x\) 处至 \(r\) 阶的导数数据，称为 \(r\) 阶\term[penshe]{喷射}[jet]。两个局部截面的这些导数全相同时代表同一喷射；链式法则保证这一等价关系与坐标选择无关。选定联络后，可以用 \(u,\nabla u,\ldots,\nabla^ru\) 的相应对称分量表示它。

\begin{proposition}\label{g3:prop:nonlinear-jet-operator}
	设 \(M^m\) 紧、光滑、无边界，\(E,F\) 为光滑丛，\(r\ge0\)、\(k>r\) 为整数，\(1<p<\infty\)、\((k-r)p>m\)。若 \(\mathcal P\) 是由至 \(r\) 阶喷射的光滑纤维映射给出的微分算子，则在其喷射值落于定义域的开集上，
	\[
		\mathcal P:W^{k,p}(M;E)\longrightarrow W^{k-r,p}(M;F)
	\]
	光滑。局部写 \(\mathcal P(u)=F(x,u,\partial u,\ldots,\partial^r u)\)，其导数为
	\[
		\mathrm D\mathcal P(u)h
		=\sum_{|\beta|\le r}\partial_{z_\beta}F(x,j^ru)\,\partial^\beta h.
	\]
\end{proposition}
\begin{proof}
	\(u\mapsto j^ru\) 在联络分解后是到有限个 \(W^{k-r,p}\) 丛截面的连续线性映射。条件 \((k-r)p>m\) 使这些分量连续，并使其空间成为乘法代数。将复合定理用于喷射纤维映射，再用 Banach 链式法则，便得结论与导数公式。有限图上的表达按照原来的微分算子变换规则粘合。
\end{proof}

这是一般光滑喷射非线性的一个充分范围。若最高阶导数仅线性出现，上一小节的乘子估计允许更弱的条件。下面的曲面方程正好处于 \(W^{1,p}\to L^p\)，其目标空间本身并非乘法代数。

\begin{theorem}\label{g3:thm:cr-nonlinear-smooth}
	设 \(\Omega\subset\mathbb R^2\) 为有界光滑域，\(p>2\)，\(O\subset\mathbb R^{2n}\) 开，\(J:O\to\operatorname{End}(\mathbb R^{2n})\) 光滑且 \(J(y)^2=-I\)。在满足连续代表 \(u(\overline\Omega)\subset O\) 的 \(W^{1,p}(\Omega;\mathbb R^{2n})\) 开集上，
	\[
		\mathcal F(u)=\partial_su+J(u)\partial_tu
	\]
	是到 \(L^p(\Omega;\mathbb R^{2n})\) 的光滑映射，且
	\begin{equation}\label{g3:eq:cr-nonlinear-derivative}
		\mathrm D\mathcal F(u)h
		=\partial_sh+J(u)\partial_th
		+\mathrm dJ(u)[h]\partial_tu.
	\end{equation}
	若 \(u\) 光滑，此线性化是光滑系数的实椭圆一阶系统。
\end{theorem}
\begin{proof}
	复合定理给出 \(u\mapsto J(u)\) 的 \(W^{1,p}\) 光滑性，而
	\(W^{1,p}\times L^p\to L^p\) 的乘法连续，因为第一因子嵌入 \(L^\infty\)。和线性导数映射复合即得光滑性。为直接看到余项，展开得
	\[
		\mathcal F(u+h)-\mathcal F(u)-\mathrm D\mathcal F(u)h
		=\bigl(J(u+h)-J(u)-\mathrm dJ(u)[h]\bigr)\partial_tu
		+\bigl(J(u+h)-J(u)\bigr)\partial_th.
	\]
	在固定紧像和有界 \(W^{1,p}\) 邻域内，其 \(L^p\) 范数至多为 \(C_R\|h\|_{W^{1,p}}^2\)。一阶导数的连续性及高阶公式也由同一复合和乘法估计得到。

	在光滑 \(u\) 处，最后一项是零阶项，主符号为
	\(\mathrm i(\sigma I+\tau J(u))\)。由 \(J^2=-I\)，其逆为
	\(-\mathrm i(\sigma I-\tau J(u))/(\sigma^2+\tau^2)\)，故在非零实余切方向可逆。
\end{proof}

这一结论可以在闭黎曼曲面 \((\Sigma,j)\) 与几乎复流形 \((N,J)\) 之间粘合。算子
\[
	\overline\partial_Ju=\frac12\bigl(\mathrm du+J(u)\mathrm du\circ j\bigr)
\]
取值于沿 \(u\) 的复反线性一形式，即满足 \(A\circ j=-J(u)\circ A\) 的实线性映射。为把变化的目标纤维变成固定 Banach 空间，在光滑参考映射 \(u_0\) 的指数图中，沿短测地线平行运输即可。可以选保持 \(J\) 的光滑联络：从任意联络出发，改为 \(\nabla^J=\nabla-\tfrac12J(\nabla J)\)，由 \(J(\nabla J)+(\nabla J)J=0\) 验证 \(\nabla^JJ=0\)。其平行运输与 \(J\) 交换，因而识别各个复反线性目标丛；系数对指数坐标光滑。上一定理于是证明该截面在固定图中为 \(W^{1,p}\to L^p\) 的光滑映射。在光滑解处，线性化适用\cref{g3:thm:elliptic-fredholm-wkp}，从而可用分裂满射或 Lyapunov--Schmidt 约化。

\subsection{椭圆自举、先验界与弱解的正则性恢复}
\label{g3:subsec:1306:05}
\begin{theorem}\label{g3:thm:semilinear-sobolev-bootstrap}
	设 \(M^m\) 为闭流形，\(P\) 为正阶 \(d\) 的光滑系数椭圆系统，\(F\) 为光滑零阶纤维映射。若 \(u\in W^{k,p}\)、\(k\ge1\)、\(1<p<\infty\)、\(kp>m\)，并在分布意义下满足 \(Pu=F(x,u)\)，则 \(u\) 光滑。在 \(\|u\|_{W^{k,p}}\) 有界且像位于固定紧集的解族中，逐次可得各固定高阶 Sobolev 范数的统一界。
\end{theorem}
\begin{proof}
	复合定理先给 \(F(\cdot,u)\in W^{k,p}\)。第十二章的分布解正则性\cref{g3:thm:elliptic-distribution-lp}在有限坐标覆盖上给出 \(u\in W^{k+d,p}\)。现在 \((k+d)p>m\)，可再次使用复合定理，得到右端属于 \(W^{k+d,p}\)，继而 \(u\in W^{k+2d,p}\)。重复即可达到任意阶，由高阶嵌入得到光滑性。每步复合估计的常数依赖前一步范数及固定紧像内有限阶系数界；全局椭圆估计又把下一阶范数控制在这些量之内，归纳即得一致性。
\end{proof}

曲面上的非线性 Cauchy--Riemann 方程，其最高阶系数 \(J(u)\) 一开始只有 \(W^{1,p}\) 正则性。对此可以直接证明以下升级，而无需先假定它已是光滑系数。

\begin{proposition}\label{g3:prop:weak-cr-smooth-bootstrap}
	设 \(\Omega\subset\mathbb R^2\)、\(p>2\)，\(J\) 在目标开集上光滑且 \(J^2=-I\)。若 \(u\in W^{1,p}_{\mathrm{loc}}(\Omega;\mathbb R^{2n})\)，连续像位于目标内，且
	\(\partial_su+J(u)\partial_tu=0\) 弱成立，则 \(u\) 光滑。
\end{proposition}
\begin{proof}
	先记录一个只需连续系数、但作用于已有 \(W^{1,q}\) 函数的估计。若 \(K(x)\) 连续、\(K^2=-I\)，\(1<q<\infty\)，则在充分小的嵌套圆盘 \(B'\Subset B\) 上有
	\begin{equation}\label{g3:eq:continuous-cr-apriori}
		\|w\|_{W^{1,q}(B')}
		\le C\bigl(\|\partial_sw+K\partial_tw\|_{L^q(B)}
		+\|w\|_{L^q(B)}\bigr),\qquad w\in W^{1,q}(B).
	\end{equation}
	证明是在中心冻结 \(K_0\)，对 \(\chi w\) 使用常系数一阶 \(L^q\) 估计。误差 \((K-K_0)\partial_t(\chi w)\) 的范数至多为 \(\|K-K_0\|_\infty\|\mathrm D(\chi w)\|_q\)，缩小球即可吸收到左侧；截断交换子为零阶，故由 \(\|w\|_q\) 控制。这个论证不微分 \(K\)。

	取 \(K=J(u)\)，它由 \(p>2\) 的嵌入而连续。在一个坐标方向取差商 \(\delta_h\)，弱方程给出
	\[
		(\partial_s+J(u)\partial_t)\delta_hu
		=-\delta_h(J(u))\,\partial_tu(\cdot+h).
	\]
	目标紧像内的链式差商满足 \(|\delta_h(J(u))|\le C|\delta_hu|\)。切片基本定理给出 \(\|\delta_hu\|_p\le C\|\mathrm Du\|_p\) 于较小区域，故右侧一致有界于 \(L^{p/2}\)。每个非零 \(h\) 的 \(\delta_hu\) 本来属于 \(W^{1,p}\)，因此可合法代入式\eqref{g3:eq:continuous-cr-apriori}，取 \(q=p/2>1\)。得到内盘上差商一致有界于 \(W^{1,p/2}\)。反射性提取弱子列，测试配对的差商极限识别它为对应的 \(\partial u\)，遂得 \(u\in W^{2,p/2}_{\mathrm{loc}}\)。

	若 \(2<p<4\)，二维嵌入使 \(\mathrm Du\in L^{p_1}_{\mathrm{loc}}\)，其中
	\[
		p_1=\frac{2p}{4-p}>p,\qquad
		\frac12-\frac1{p_1}=2\left(\frac12-\frac1p\right).
	\]
	重复有限次便达到 \(p\ge4\)；等号时的临界嵌入给出任意有限的更高指数。因此最终 \(u\in W^{2,q}_{\mathrm{loc}}\) 对某个 \(q>2\) 成立，进而 \(u\in C^{1,\alpha}_{\mathrm{loc}}\)，某个 \(\alpha>0\)。每步只需再缩小内部区域。

	现在可在弱意义下对原方程作用 \(\partial_s-J(u)\partial_t\)，用乘积法则和 \(J^2=-I\) 得
	\[
		(\partial_s^2+\partial_t^2)u
		=-\bigl(\partial_sJ(u)-J(u)\partial_tJ(u)\bigr)\partial_tu.
	\]
	右侧属于 \(C^{0,\alpha}_{\mathrm{loc}}\)。此时主部已是光滑的标量 Laplace 算子，可逐分量应用分布型 Schauder 正则性\cref{g3:cor:schauder-distribution-smooth-coefficients}，得到 \(u\in C^{2,\alpha}_{\mathrm{loc}}\)。右端随之提升一阶，再反复应用同一结论，即得光滑性。
\end{proof}

对一般准线性系统，起始正则性、主部对 \(u\) 及其导数的依赖，以及冻结后可吸收的范数，都须像上述证明一样逐项确定。先验界则承担另一职责：它使一个变化的解族保持在同一组估计常数所覆盖的区域；单个解光滑并不直接给出这样的统一界。

\subsection{直接方法、约束临界点与规范对称的区别}
\label{g3:subsec:1306:06}
局部隐函数方法从一个已知解及其线性化出发。直接方法从极小化序列出发；约束改变允许的变分方向，规范对称则使多个未知量代表同一几何对象。以下模型分别落实这三种结构。

\begin{example}\label{g3:exm:semilinear-convex-minimization}
	在闭黎曼曲面 \(M\) 上，给定实光滑 \(f\)，考虑
	\[
		E(u)=\frac12\int_M(|\mathrm du|^2+u^2)\,\mu_g
		+\frac14\int_Mu^4\,\mu_g-\int_Mfu\,\mu_g,
		\qquad u\in H^1(M;\mathbb R).
	\]
	二维嵌入 \(H^1\to L^4\) 使它有定义。由 Cauchy--Schwarz 与平方配方，\(E(u)\ge\tfrac14\|u\|_{H^1}^2-C_f\)，所以极小化序列有界。弱子列在 \(H^1\) 中收敛，Rellich 紧性使它在 \(L^4\) 中强收敛；二次范数弱下半连续，四次项与线性项收敛，故极小值取得。严格凸性给出唯一性。

	对任意 \(h\in H^1\) 微分，得到弱方程
	\[
		\Delta_gu+u+u^3=f.
	\]
	由于 \(H^1\to L^6\)，右端 \(f-u-u^3\in L^2\)，椭圆正则性先给 \(u\in H^2\)；此时 \(2\cdot2>\dim M\)，上一半线性自举定理给出 \(u\) 光滑。这一存在证明使用能量与紧性，并不要求事先知道一个解。
\end{example}

\begin{proposition}\label{g3:prop:constraint-lagrange-banach}
	设 \(X\) 为实 Banach 空间，\(E:X\to\mathbb R\)、\(C:X\to\mathbb R^a\) 在 \(u\) 附近为 \(C^1\)，且 \(\mathrm DC(u)\) 满射。若 \(u\) 为约束 \(C=C(u)\) 上的局部极值，则存在唯一 \(\lambda\in(\mathbb R^a)^*\)，使
	\[
		\mathrm DE(u)=\lambda\circ\mathrm DC(u).
	\]
\end{proposition}
\begin{proof}
	有限维目标的满射有有界右逆：逐个提升一组基并线性延拓即可。分裂水平集定理说明每个 \(h\in\ker\mathrm DC(u)\) 都是约束中一条 \(C^1\) 曲线的初速度，故 \(\mathrm DE(u)h=0\)。取右逆 \(Q\)，定义 \(\lambda=\mathrm DE(u)\circ Q\)。分解
	\(h=(h-Q\mathrm DC(u)h)+Q\mathrm DC(u)h\)，第一项在核中，遂得公式。满射性保证 \(\lambda\) 唯一。
\end{proof}

例如将二次能量 \(\tfrac12\int(|\mathrm du|^2+Vu^2)\) 限制在 \(\int u^2=1\) 上，得到的临界方程是 \(\Delta_gu+Vu=\lambda u\)，其中乘子来自约束，而非原先的无约束 Euler 方程。直接方法还须证明约束在所用极限下保留；这里通过 \(H^1\) 到 \(L^2\) 的紧性保持单位范数，不能仅凭弱收敛断言。

\ProseParagraphBreak

规范对称可在一形式上直接看到。令 \(a\) 为闭连通流形上的光滑实一形式，能量 \(\tfrac12\|\mathrm da\|_2^2\) 在 \(a\mapsto a+\mathrm d\phi\) 下不变，所以 Euler 算子 \(\mathrm d^*\mathrm d\) 的核包含所有恰当一形式。它的主符号也沿 \(\xi\) 方向退化。选取代表条件 \(\mathrm d^*a=0\)，需解
\[
	\Delta_g\phi=-\mathrm d^*a.
\]
右端与常数正交，因为 \(\langle\mathrm d^*a,1\rangle=\langle a,\mathrm d1\rangle=0\)，故零均值的 \(\phi\) 存在且唯一。此时 \(a+\mathrm d\phi\) 满足代表条件，而主方程与该条件组合成 \((\mathrm d^*\mathrm d+\mathrm d\mathrm d^*)a\) 的椭圆结构。这里解的是如何在同一规范类中选代表；剩余调和形式由前面的椭圆复形理论控制。对于度量、自微分同胚或非交换规范群，代表条件、稳定子及作用的微分损失需要按相应对象重新建立。
% 纲要条目排为编号小节标题，并进入目录与书签；说明段仍保留于纲要注释中。


\begin{chaptersummary}
闭流形上的局部估计经有限覆盖成为全局先验界，Rellich 紧性使椭圆核有限维，并给出去掉核后的加强估计。由此，收敛的右端产生收敛的归一化解，值域闭；与值域正交的截面恰是伴随齐次解。Fredholm 选择将可解性写成有限个积分条件，同一光滑核与余核贯穿各个整数阶、\(1<p<\infty\) 的 Sobolev 实现。指数在连续扰动下稳定，核维数及逆的范数则可以变化。

\ProseParagraphBreak

Green 算子选择核正交的规范解，并以有限秩投影记录求逆误差。非负自伴算子的一个预解式是紧正算子，二次型最大化由此产生完备的光滑特征基；带权平方和记录算子的整数次幂和相应 Sobolev 正则性。椭圆复形的符号正合保证其 Laplace 算子椭圆，Green 交换关系进而给出正合、余正合与调和正交分解，每个光滑上同调类具有唯一调和代表。

\ProseParagraphBreak

非线性求解把线性估计用于真实余项。定量隐函数方法控制解的邻域与参数依赖，分裂满射产生局部水平集，有限维核与余核则留下障碍映射。Sobolev 连续嵌入和乘法估计使映射坐标、复合及线性化成为有界且可微的操作；微分算子的目标空间须记录导数损失。最后，正则性自举、变分极小化和对称性处理分别承担提升解、取得解和选择代表的任务，应用时需要逐项匹配。
\end{chaptersummary}

\BookExercises

\begin{exercise}\label{g3:ex:periodic-drift-fredholm}
	设 \(a,f\in C^\infty(S^1;\mathbb R)\)，\(S^1=\mathbb R/(2\pi\mathbb Z)\)。记 \(A(x)=\int_0^x a(t)\,\mathrm dt\)、\(\tau=A(2\pi)\)。求周期方程 \(u'+au=f\) 的全部解，并分别在 \(\tau=0\) 与 \(\tau\ne0\) 时计算 \(P=\partial_x+a\) 及其形式伴随的核。验证所得可解条件与 Fredholm 选择相符。
\end{exercise}
\begin{solution}
	乘以积分因子，得到 \((e^Au)'=e^Af\)，所以
	\[
		u(x)=e^{-A(x)}\left(c+\int_0^x e^{A(t)}f(t)\,\mathrm dt\right).
	\]
	周期条件等价于
	\((e^\tau-1)c=\int_0^{2\pi}e^{A(t)}f(t)\,\mathrm dt\)。若 \(\tau\ne0\)，它唯一确定 \(c\)，故任意光滑右端有唯一解。若 \(\tau=0\)，相容条件为该积分等于零，此时 \(c\) 任意。端点值相等后，方程及周期系数逐阶确定相同的端点导数，故所得解确为光滑周期函数。

	分部积分给出 \(P^*=-\partial_x+a\)。两齐次方程的局部解分别为 \(c e^{-A}\) 与 \(c e^A\)。当 \(\tau\ne0\) 时均无非零周期解；当 \(\tau=0\) 时，它们各生成一个一维核。此时相容条件正是 \((f,e^A)_2=0\)。两个核通常由不同函数生成，不能把“与伴随核正交”改成“与原核正交”；两种情形的指数都为零。
\end{solution}

\begin{exercise}\label{g3:ex:noncompact-derivative-range}
	在 \(\mathbb R\) 上考虑 \(P=\partial_x:H^1(\mathbb R)\to L^2(\mathbb R)\)。证明其主符号椭圆、核为零，但值域不闭。具体取 \(\chi\in C_c^\infty(\mathbb R)\) 在零点附近为一，定义
	\[
		\widehat f=\chi,\qquad
		\widehat f_j=\mathbf1_{\{|\xi|\ge1/j\}}\chi.
	\]
	证明 \(f_j\) 在值域中、\(f_j\to f\) 于 \(L^2\)，而 \(f\) 不在值域中。再取非零 \(\eta\in C_c^\infty(\mathbb R)\)，用 \(u_R(x)=R^{-1/2}\eta(x/R)\) 说明补空间估计缺失的机制。
\end{exercise}
\begin{solution}
	符号为 \(\mathrm i\xi\)，在实 \(\xi\ne0\) 时非零。若 \(u'=0\) 分布成立，\(u\) 为常数；属于 \(L^2(\mathbb R)\) 的常数只有零。也可由 \(\xi\widehat u=0\) 且 \(\widehat u\in L^2\) 立即得到核为零。

	令 \(\widehat v_j=\widehat f_j/(\mathrm i\xi)\)，并在 \(\xi=0\) 取零。每个 \(j\) 的频率支集紧且避开零点，故 \(v_j\in H^1\)，且 \(v_j'=f_j\)。Plancherel 定理与控制收敛给出 \(f_j\to f\)。若 \(u'=f\)、\(u\in H^1\)，则在 \(\xi\ne0\) 几乎处处有 \(\widehat u=\chi/(\mathrm i\xi)\)，但它在零点附近不平方可积，矛盾。

	最后，\(\|u_R\|_2=\|\eta\|_2\)，而 \(\|u_R'\|_2=R^{-1}\|\eta'\|_2\to0\)。没有常数 \(C\) 能使 \(\|u\|_{H^1}\le C\|u'\|_2\)。这些截面的质量在越来越大的尺度上分散；闭流形证明使用的全局 Rellich 紧性在此不成立。局部椭圆性本身并未失败。
\end{solution}

\begin{exercise}\label{g3:ex:shift-fredholm-index}
	在 \(\ell^2(\mathbb N;\mathbb C)\) 上定义
	\[
		R(x_1,x_2,\ldots)=(0,x_1,x_2,\ldots),\qquad
		L(x_1,x_2,\ldots)=(x_2,x_3,\ldots).
	\]
	计算它们的伴随、核、余核与指数。对 \(\alpha\in\mathbb C\)，再考虑 \(R_\alpha x=Rx+\alpha x_1e_1\)，直接证明其 Fredholm 指数独立于 \(\alpha\)，并写出其值域的一个方程。
\end{exercise}
\begin{solution}
	直接逐项配对得 \(R^*=L\)、\(L^*=R\)。\(R\) 是等距嵌入，核为零，值域是第一坐标为零的闭超平面，故余核一维、指数为 \(-1\)。\(L\) 满射，核为 \(\mathbb C e_1\)，故余核零维、指数为一。

	有 \(R_\alpha x=(\alpha x_1,x_1,x_2,\ldots)\)，所以它仍单射，且
	\[
		\operatorname{ran}R_\alpha
		=\{y\in\ell^2:y_1=\alpha y_2\}.
	\]
	反向把满足条件的 \(y\) 送到 \((y_2,y_3,\ldots)\)，即可验证值域等式。右边是连续非零线性泛函 \(y\mapsto y_1-\alpha y_2\) 的核，闭且余维一，所以指数始终为 \(-1\)。这个例子说明一般 Fredholm 算子可以单射而不满射；“单射等价于满射”还需要指数为零。
\end{solution}

\begin{exercise}\label{g3:ex:rank-one-resonance}
	设 \(M\) 为非空正维闭黎曼流形，\(h\in C^\infty(M;\mathbb R)\) 非零，令 \(A=I+\Delta_g\)、\(w=A^{-1}h\) 和 \(a=(w,h)_2\)。证明 \(a>0\)。对实数 \(t\)，研究
	\[
		A u-t(u,h)_2h=f,\qquad f\in L^2(M;\mathbb R),\quad u\in H^2(M).
	\]
	给出 \(ta\ne1\) 时的唯一解，以及 \(ta=1\) 时的核与相容条件。证明这族有界映射 \(H^2\to L^2\) 始终有指数零，并解释逆在共振参数附近的变化。
\end{exercise}
\begin{solution}
	\(A\) 形式自伴且能量为 \(\|u\|_2^2+\|\mathrm du\|_2^2\)，故核为零。Fredholm 选择给出有界逆，正则性使 \(w\) 光滑，且
	\[
		a=(w,Aw)_2=\|w\|_2^2+\|\mathrm dw\|_2^2>0.
	\]
	记 \(v=A^{-1}f\)。原方程等价于 \(u=v+t(u,h)_2w\)，与 \(h\) 配对得到
	\((1-ta)(u,h)_2=(v,h)_2\)。因此 \(ta\ne1\) 时唯一解为
	\[
		u=v+\frac{t(v,h)_2}{1-ta}\,w.
	\]
	当 \(ta=1\) 时，必要条件为 \((v,h)_2=0\)。由自伴性，\((v,h)_2=(v,Aw)_2=(Av,w)_2=(f,w)_2\)。这个条件成立时，全部解为 \(v+c w\)，\(c\in\mathbb R\)；其核就是 \(\mathbb Rw\)，值域为闭超平面 \(w^\perp\)。因此非共振参数处核与余核都为零，共振参数处二者都一维，指数始终为零。

	取 \(f=h\)，有 \(v=w\)、\((v,h)_2=a\)，解变为 \(u=w/(1-ta)\)。其范数在 \(t\to1/a\) 时发散，给出逆算子不能一致有界的直接证据。有限秩项不是局部微分项，但其全部可解性仍可由这个代数计算确定。
\end{solution}

\begin{exercise}\label{g3:ex:boundary-free-infinite-kernel}
	在单位圆盘 \(D\subset\mathbb R^2\) 上，将 \(\Delta_0=-\partial_x^2-\partial_y^2\) 看成 \(H^2(D)\to L^2(D)\) 的有界映射，不规定边界条件。证明它的核无限维，因而不是 Fredholm 算子。再说明要求 \(u\in H_0^1(D)\cap H^2(D)\) 如何消除这个核，并指出这个计算还没有单独证明什么。
\end{exercise}
\begin{solution}
	对每个 \(j\ge0\)，\(u_j(x,y)=\operatorname{Re}(x+\mathrm iy)^j\) 是光滑调和多项式，属于 \(H^2(D)\)。它们线性无关：若有限线性组合为零，限制到实轴小区间就得到一个恒零多项式，所有系数均为零。因此未指定边界条件时核无限维，不满足 Fredholm 定义。

	若 \(u\in H_0^1(D)\cap H^2(D)\) 且 \(\Delta_0u=0\)，以光滑紧支函数在 \(H_0^1\) 中逼近 \(u\)，弱分部积分给出 \(\int_D|\mathrm Du|^2=0\)。零边界 Poincaré 不等式使 \(u=0\)。这证明了 Dirichlet 问题的唯一性；对每个 \(L^2\) 右端的整体 \(H^2\) 存在性还需要边界正则性，不能仅由核为零代替。
\end{solution}

\begin{exercise}\label{g3:ex:elliptic-variational-eigenvalues}
	设 \(M\) 是非空闭黎曼流形，\(\dim M\ge1\)，\(E\to M\) 是正秩复 Hermitian 丛，\(P\) 是光滑、形式自伴、非负的 \(d\ge1\) 阶椭圆算子。将特征值按重数写作 \(\lambda_1\le\lambda_2\le\cdots\)。证明
	\[
		\lambda_j
		=\inf_{\substack{V\subset H^d(E)\\\dim V=j}}
		\sup_{0\ne u\in V}\frac{\langle Pu,u\rangle}{\|u\|_2^2}.
	\]
	再设 \(W\) 为光滑 Hermitian 丛自同态场，逐点满足 \(0\le W\le cI\)，其中 \(c\ge0\) 为常数。证明
	\[
		\lambda_j(P)\le\lambda_j(P+W)\le\lambda_j(P)+c
		\quad(j\ge1).
	\]
\end{exercise}
\begin{solution}
	设 \((e_k)\) 是正交归一特征基。对 \(u\in H^d\)，\cref{g3:thm:elliptic-discrete-spectrum}给出
	\[
		\langle Pu,u\rangle=\sum_k\lambda_k|\langle u,e_k\rangle|^2.
	\]
	级数绝对收敛，因为 \(u\in L^2\)、\(Pu\in L^2\)，可用 Cauchy--Schwarz 控制 \(\sum_k\lambda_k|\langle u,e_k\rangle|^2\)。

	任取 \(j\) 维子空间 \(V\subset H^d\)。将 \(u\in V\) 映到前 \(j-1\) 个系数，是从 \(j\) 维空间到 \(\mathbb C^{j-1}\) 的线性映射，因此有非零核向量 \(u\)。它的前 \(j-1\) 个系数为零，故
	\[
		\langle Pu,u\rangle
		=\sum_{k\ge j}\lambda_k|\langle u,e_k\rangle|^2
		\ge\lambda_j\|u\|_2^2.
	\]
	所以每个 \(V\) 上的上确界至少为 \(\lambda_j\)。另一方面，取 \(V=\operatorname{span}(e_1,\ldots,e_j)\)，同一展开使商不超过 \(\lambda_j\)，而取 \(u=e_j\) 达到它。这证明所求公式，且此处下确界实际能够达到。

	\(P+W\) 的最高阶符号与 \(P\) 相同，仍椭圆、形式自伴且非负，其 \(L^2\) 实现域仍为 \(H^d\)。逐点 Hermitian 不等式积分后给出，对每个非零 \(u\in H^d\)，
	\[
		\frac{\langle Pu,u\rangle}{\|u\|_2^2}
		\le\frac{\langle(P+W)u,u\rangle}{\|u\|_2^2}
		\le\frac{\langle Pu,u\rangle}{\|u\|_2^2}+c.
	\]
	先在每个 \(j\) 维子空间上取上确界，再对全部这样的子空间取下确界，即得两个特征值界。估计不要求两个算子的特征基相同。
\end{solution}

\begin{exercise}\label{g3:ex:green-resolvent-regularization}
	设 \(M\) 为非空、紧且无边界的正维光滑黎曼流形，\(E\to M\) 为正秩光滑复 Hermitian 丛，\(P\) 为光滑、形式自伴、非负的 \(d\ge1\) 阶椭圆算子。令 \(\Pi\) 为其零空间的 \(L^2\) 正交投影，\(G\) 为 Green 算子，\(\lambda_+>0\) 为最小正特征值。对 \(\varepsilon>0\)，令 \(R_\varepsilon=(P+\varepsilon I)^{-1}\)，其逆按 \(H^d\to L^2\) 理解。证明
	\[
		\left\|R_\varepsilon-\varepsilon^{-1}\Pi-G\right\|_{L^2\to L^2}
		\le\frac{\varepsilon}{\lambda_+^2}.
	\]
	再证明 \(\varepsilon R_\varepsilon\to\Pi\) 按 \(L^2\) 算子范数收敛，并说明 \(R_\varepsilon\) 在 \(\varepsilon\downarrow0\) 时有界的充要条件。
\end{exercise}
\begin{solution}
	所有 \(\lambda_j\ge0\)，故 \(-\varepsilon\) 在预解集中。逐个特征向量上有
	\[
		R_\varepsilon e_j=(\lambda_j+\varepsilon)^{-1}e_j.
	\]
	当 \(\lambda_j=0\) 时，\(R_\varepsilon-\varepsilon^{-1}\Pi-G\) 的系数为零；当 \(\lambda_j>0\) 时，其系数为
	\[
		\frac1{\lambda_j+\varepsilon}-\frac1{\lambda_j}
		=-\frac{\varepsilon}{\lambda_j(\lambda_j+\varepsilon)}.
	\]
	绝对值至多 \(\varepsilon/\lambda_+^2\)。对任意 \(L^2\) 展开求平方和，便得所述算子范数估计。

	同理，\(\varepsilon R_\varepsilon-\Pi\) 在零空间为零，在正特征向量上乘以 \(\varepsilon/(\lambda_j+\varepsilon)\)，所以
	\[
		\|\varepsilon R_\varepsilon-\Pi\|_{L^2\to L^2}
		=\frac{\varepsilon}{\lambda_++\varepsilon}\longrightarrow0.
	\]
	若零空间非零，在其中取单位向量 \(h\)，则 \(\|R_\varepsilon h\|_2=\varepsilon^{-1}\)，因而无一致界。若零空间为零，则 \(\lambda_1=\lambda_+>0\)，并有 \(\|R_\varepsilon\|=(\lambda_1+\varepsilon)^{-1}\le\lambda_1^{-1}\)。所以有界的充要条件为 \(\ker P=0\)。发散部分恰好是有限秩项 \(\varepsilon^{-1}\Pi\)，将它扣除才得到 Green 算子的极限。
\end{solution}

\begin{exercise}\label{g3:ex:circle-green-kernel}
	在 \(S^1=\mathbb R/(2\pi\mathbb Z)\) 上，取非负 Laplace 算子 \(P=-\partial_x^2\)。将函数
	\[
		K(t)=\frac{t^2}{4\pi}-\frac{|t|}2+\frac\pi6,
		\qquad -\pi\le t\le\pi,
	\]
	按 \(2\pi\) 周期延拓。证明 \(K\) 连续且均值为零，并在周期分布意义下满足
	\[
		-K''=\delta_0-\frac1{2\pi}.
	\]
	证明对每个 \(f\in L^2(S^1)\)，其规范 Green 解可写成
	\[
		Gf(x)=\int_{-\pi}^\pi K(x-y)f(y)\,\mathrm dy.
	\]
	最后由特征展开证明
	\(K(t)=\frac1{2\pi}\sum_{k\ne0}e^{\mathrm ikt}/k^2\)，并从 \(t=0\) 求出 \(\sum_{k\ge1}k^{-2}\)。
\end{exercise}
\begin{solution}
	两端函数值同为 \(-\pi/12\)，而两端的一阶导数都为零，所以周期延拓在识别端点处没有额外的导数跳跃。零点两侧函数连续，一阶导数分别为 \(t/(2\pi)+1/2\) 与 \(t/(2\pi)-1/2\)，在零点有大小 \(-1\) 的跳跃。因此弱求导给出 \(K''=1/(2\pi)-\delta_0\)。直接积分为
	\[
		\int_{-\pi}^\pi K(t)\,\mathrm dt
		=2\left(\frac{\pi^2}{12}-\frac{\pi^2}4+\frac{\pi^2}6\right)=0.
	\]
	连续性已经由分段公式及两处接合检查得到。

	令 \(u\) 为题中的卷积。\(K\) 有界而圆周体积有限，由 Cauchy--Schwarz 得 \(u\in L^2\)，并且积分换序给出 \(\int u=0\)。对光滑周期测试函数交换积分，使用已证的分布恒等式，即得
	\[
		-u''=f-\frac1{2\pi}\int_{-\pi}^\pi f.
	\]
	右端在 \(L^2\)，椭圆正则性给出 \(u\in H^2\)；零均值条件使它正交于常数核。因此由\cref{g3:thm:elliptic-green-general}的唯一性，\(u=Gf\)。这个卷积公式中的积分使用通常长度测度，故不再补乘 \(1/(2\pi)\)。

	对分布方程取第 \(k\) 个周期 Fourier 系数。采用 \(\widehat K(k)=(2\pi)^{-1}\int K(t)e^{-\mathrm ikt}\,\mathrm dt\)，零均值给出 \(\widehat K(0)=0\)，而 \(k\ne0\) 时
	\(k^2\widehat K(k)=1/(2\pi)\)。由圆周特征基的完备性，所给级数在 \(L^2\) 中等于 \(K\)。因 \(\sum_{k\ne0}k^{-2}<\infty\)，该级数还绝对一致收敛到连续函数；两个连续代表既几乎处处相等，就处处相等。取 \(t=0\)，得到
	\[
		\frac\pi6=\frac1\pi\sum_{k=1}^\infty\frac1{k^2},
		\qquad \sum_{k=1}^\infty\frac1{k^2}=\frac{\pi^2}6.
	\]
\end{solution}

\begin{exercise}\label{g3:ex:twisted-circle-complex}
	在 \(S^1=\mathbb R/2\pi\mathbb Z\) 上，对实数 \(a\) 定义两项复形
	\[
		C^\infty(S^1;\mathbb C)\xrightarrow{D_a}
		C^\infty(S^1;\mathbb C),\qquad D_af=f'+\mathrm i a f.
	\]
	证明它椭圆，计算 \(H_{D_a}^0,H_{D_a}^1\) 及两个次数的调和空间。给出方程 \(D_af=g\) 的可解条件，并在 \(a\notin\mathbb Z\) 时写出唯一周期解。解释核维数的变化为何不改变该两项复形的指数。
\end{exercise}
\begin{solution}
	主符号为 \(\mathrm i\xi\)，与 \(a\) 无关，在 \(\xi\ne0\) 时可逆。齐次解为 \(f(x)=Ce^{-\mathrm iax}\)，周期性要求
	\((e^{-2\pi\mathrm ia}-1)C=0\)。因此当 \(a\in\mathbb Z\) 时，核是一维空间 \(\mathbb C e^{-\mathrm iax}\)；其余参数处核为零。

	乘以积分因子 \(e^{\mathrm iax}\)，非齐次方程变成 \((e^{\mathrm iax}f)'=e^{\mathrm iax}g\)。所有解为
	\[
		f(x)=e^{-\mathrm iax}\left(C+\int_0^x e^{\mathrm iat}g(t)\,\mathrm dt\right).
	\]
	若 \(a\notin\mathbb Z\)，唯一使其周期的常数为
	\[
		C=\frac{\int_0^{2\pi}e^{\mathrm iat}g(t)\,\mathrm dt}
		{e^{2\pi\mathrm ia}-1}.
	\]
	端点处函数值匹配；由方程及 \(g\) 的周期性，依次求导可知所有导数也匹配，因而确实定义圆周上的光滑解。此时像为整个靶空间，两个上同调都为零。

	若 \(a\in\mathbb Z\)，可解条件恰为
	\(\int_0^{2\pi}e^{\mathrm iat}g(t)\,\mathrm dt=0\)，其成立时 \(C\) 任意。这个条件是一个非零线性泛函的零集，所以余核一维。对通常 \(L^2\) 内积，\(D_a^*=-\partial_x-\mathrm ia=-D_a\)，故两个次数的 Laplace 算子均为 \(-D_a^2\)，调和空间都等于 \(\ker D_a\)。在整数参数处，核与余核同时增加一维，故指数 \(\dim H^0-\dim H^1\) 始终为零。
\end{solution}

\begin{exercise}\label{g3:ex:flat-torus-explicit-hodge}
	在平坦环面 \(\mathbb T^2=(\mathbb R/2\pi\mathbb Z)^2\) 上，令
	\[
		\alpha=(a+\cos x)\,\mathrm dx+(b+\sin x)\,\mathrm dy,
		\qquad a,b\in\mathbb R.
	\]
	求 \(\alpha\) 的调和、正合和余正合三个分量。分别求可选择的零均值函数原像和零均值二形式原像，并计算三分量的 \(L^2\) 范数平方。判断 \(\alpha\) 能否代表 de Rham 上同调类。
\end{exercise}
\begin{solution}
	分部积分或外乘收缩公式给出
	\[
		\mathrm d^*(h\,\mathrm dx\wedge\mathrm dy)
		=(\partial_yh)\,\mathrm dx-(\partial_xh)\,\mathrm dy.
	\]
	因此
	\[
		\alpha=\underbrace{a\,\mathrm dx+b\,\mathrm dy}_{\text{调和}}
		+\underbrace{\mathrm d(\sin x)}_{\text{正合}}
		+\underbrace{\mathrm d^*(\cos x\,\mathrm dx\wedge\mathrm dy)}_{\text{余正合}}.
	\]
	两个所写原像的系数均零均值。常系数形式调和，而后三项处于相应微分像中，故由正文的正交唯一性，这正是所求分解。直接积分得
	\[
		\|a\,\mathrm dx+b\,\mathrm dy\|_2^2=4\pi^2(a^2+b^2),
		\qquad
		\|\cos x\,\mathrm dx\|_2^2
		=\|\sin x\,\mathrm dy\|_2^2=2\pi^2.
	\]
	另外，\(\mathrm d\alpha=\cos x\,\mathrm dx\wedge\mathrm dy\ne0\)，所以 \(\alpha\) 不是闭形式，不能直接代表上同调类。调和投影对所有形式都有定义，而“同一个上同调类的调和代表”只对闭形式成立。
\end{solution}

\begin{exercise}\label{g3:ex:exact-one-form-twisted-complex}
	在平坦 \(\mathbb T^2\) 上取光滑实函数 \(\phi\)，在复值微分形式上定义 \(D_\phi=\mathrm d+\mathrm d\phi\wedge\)。证明 \(D_\phi^2=0\) 且它是椭圆复形，并证明乘以 \(e^\phi\) 诱导其上同调与普通 de Rham 上同调的同构。使用通常的平坦 \(L^2\) 内积，求它在零次与二次的调和截面；当 \(\phi\) 不恒定时，解释同构的上同调为何可以有不同形状的调和代表。
\end{exercise}
\begin{solution}
	对任意形式 \(u\)，乘积法则给出
	\[
		D_\phi u=e^{-\phi}\mathrm d(e^\phi u).
	\]
	因此 \(D_\phi^2=e^{-\phi}\mathrm d^2e^\phi=0\)，且乘以 \(e^\phi\) 是与微分交换的可逆映射，它将闭形式与正合形式分别双射到普通的闭形式与正合形式，因而诱导上同调同构。\(\mathrm d\phi\wedge\) 是零阶算子，所以 \(D_\phi\) 与 \(\mathrm d\) 有相同主符号，复形椭圆。

	零次调和条件只有 \(D_\phi f=0\)，即 \(\mathrm d(e^\phi f)=0\)。环面连通，故所有零次调和截面为 \(Ce^{-\phi}\)。二次形式写作 \(h\,\mathrm dx\wedge\mathrm dy\)，此时向前的微分自动为零，只需计算伴随条件。外乘实一形式的伴随为相应向量的收缩，故
	\[
		D_\phi^*=\mathrm d^*+\iota_{\operatorname{grad}\phi},
	\]
	并且
	\[
		D_\phi^*(h\,\mathrm dx\wedge\mathrm dy)
		=(h_y-\phi_yh)\,\mathrm dx+(-h_x+\phi_xh)\,\mathrm dy.
	\]
	其为零等价于 \(\mathrm d(e^{-\phi}h)=0\)，所以二次调和空间为 \(\mathbb C e^\phi\,\mathrm dx\wedge\mathrm dy\)。这两个空间都仍是一维，然而代表一般不再为常系数。乘以 \(e^\phi\) 保持复形关系，却不是通常 \(L^2\) 内积的等距映射；上同调同构因此不要求保持指定度量选出的调和代表。
\end{solution}

\begin{exercise}\label{g3:ex:elliptic-complex-euler-index}
	设 \((E^\bullet,D)\) 为正维紧致无边界光滑流形上的一阶椭圆复形，固定光滑 Hermitian 度量与正密度。令
	\[
		E^{\mathrm{ev}}=\bigoplus_{q\ \mathrm{even}}E^q,
		\qquad E^{\mathrm{odd}}=\bigoplus_{q\ \mathrm{odd}}E^q,
		\qquad \mathcal Q=D+D^*.
	\]
	证明 \(\mathcal Q^+:C^\infty(E^{\mathrm{ev}})\to C^\infty(E^{\mathrm{odd}})\) 为椭圆算子，并证明
	\[
		\operatorname{ind}\mathcal Q^+
		=\sum_q(-1)^q\dim H_D^q(M).
	\]
	说明这个等式中的解析核、余核分别对应哪些次数的调和截面。
\end{exercise}
\begin{solution}
	\(D\) 提高一次，\(D^*\) 降低一次，所以 \(\mathcal Q\) 交换偶次与奇次部分。由于 \(D^2=(D^*)^2=0\)，有
	\[
		\mathcal Q^2=DD^*+D^*D=\bigoplus_q\Delta_q.
	\]
	在非零余切方向上，右侧符号正定可逆。于是 \(\sigma_1(\mathcal Q)\) 可逆；它的两个非对角块互为伴随，并分别给出偶纤维到奇纤维、奇纤维到偶纤维的双射。特别地两丛秩相同，\(\mathcal Q^+\) 确为方形椭圆算子，其形式伴随是另一块 \(\mathcal Q^-\)。

	对光滑总截面 \(u=\sum_qu_q\)，形式自伴与分次正交性给出
	\[
		\|\mathcal Qu\|_2^2
		=(\mathcal Q^2u,u)
		=\sum_q\bigl(\|D_qu_q\|_2^2+\|D_{q-1}^*u_q\|_2^2\bigr).
	\]
	因此 \(\mathcal Q u=0\) 当且仅当每个 \(u_q\) 调和。椭圆正则性使 Sobolev 核同样由这些光滑截面组成，故 \(\ker\mathcal Q^+=\bigoplus_{q\ \mathrm{even}}\mathcal H_D^q\)，而伴随核为奇次调和空间之和。\cref{g3:thm:elliptic-fredholm}将余核识别为这个伴随核；再用调和代表定理 \(\dim\mathcal H_D^q=\dim H_D^q(M)\)，便得到所述指数公式。
\end{solution}

\begin{exercise}\label{g3:ex:directional-not-frechet}
	在 \(\mathbb R^2\) 上令
	\[
		f(x,y)=\begin{cases}\displaystyle\frac{x^3y}{x^4+y^2},&(x,y)\ne(0,0),\\0,&(x,y)=(0,0).\end{cases}
	\]
	证明它在原点连续，沿每个固定方向的方向导数都为零，但在原点不 Fréchet 可微。指出哪个估计区分了这两种可微性。
\end{exercise}
\begin{solution}
	由 \(2x^2|y|\le x^4+y^2\)，有 \(|f(x,y)|\le|x|/2\)，所以它在原点连续。沿方向 \((a,b)\)，若 \(b\ne0\)，则
	\[
		\frac{f(ta,tb)}t=\frac{t a^3b}{t^2a^4+b^2}\longrightarrow0;
	\]
	若 \(b=0\)，分子恒为零。因此所有方向导数存在，合起来也是零线性映射。

	若存在 Fréchet 微分，它必须等于这个零映射。但是沿 \(y=x^2\) 有 \(f(x,x^2)=x/2\)，故
	\[
		\frac{|f(x,x^2)|}{\sqrt{x^2+x^4}}
		=\frac1{2\sqrt{1+x^2}}\longrightarrow\frac12.
	\]
	余项与增量范数的比值并不趋零。方向导数只固定一条直线，本例中破坏统一估计的增量方向随 \(x\) 改变。
\end{solution}

\begin{exercise}\label{g3:ex:lyapunov-schmidt-coupled-branch}
	设 \(H\) 为实 Hilbert 空间，\(w\in H\)、\(\ell\in H^*\)，记 \(a=\ell(w)\)。考虑
	\[
		F_\lambda(t,z)=(t^3-\lambda t+\ell(z),\ z-t^2w),
		\qquad(t,z)\in\mathbb R\oplus H.
	\]
	对 \(F_0\) 在原点的线性化求核、像与余核，写出一个 Lyapunov--Schmidt 障碍。分别在 \(a=0\) 与 \(a\ne0\) 时，描述 \(\lambda=0\) 附近经过原点的实解分支。
\end{exercise}
\begin{solution}
	线性化为 \(L(t,z)=(\ell(z),z)\)。其核是 \(\mathbb R\oplus\{0\}\)，像是闭图 \(\{(\ell(z),z):z\in H\}\)，可取余核补空间 \(\mathbb R\oplus\{0\}\)。值域投影取 \(\Pi_R(b,y)=(\ell(y),y)\)，余核投影为 \(\Pi_C(b,y)=(b-\ell(y),0)\)。值域方程等价于 \(z=t^2w\)，所以障碍为
	\[
		\kappa_\lambda(t)=t^3+a t^2-\lambda t
		=t(t^2+a t-\lambda).
	\]
	因此始终有平凡分支 \((t,z)=(0,0)\)。若 \(a=0\)，另外两个实分支在 \(\lambda>0\) 时为
	\(t=\pm\sqrt\lambda\)、\(z=\lambda w\)，\(\lambda<0\) 时没有非零实解。

	若 \(a\ne0\)，\(t^2+at-\lambda=0\) 在 \((0,0)\) 的 \(t\) 导数为 \(a\)，故有唯一的光滑小分支 \(t=\tau(\lambda)\)，满足 \(\tau(0)=0\)、\(\tau'(0)=1/a\)，相应 \(z=\tau(\lambda)^2w\)。另一个根在 \(\lambda=0\) 时为 \(-a\)，不经过原点，可用小邻域排除。因此非零的二次障碍系数改变了分支形状，尽管两种情形的线性化核与余核维数相同。
\end{solution}

\begin{exercise}\label{g3:ex:semilinear-cubic-expansion}
	设 \(M^m\) 为闭黎曼流形，整数 \(s>m/2\)，令 \(A=I+\Delta_g\)。证明在 \(f=0\) 附近，方程
	\[
		Au+u^3=f
	\]
	给出唯一的充分小解 \(u\in H^{s+2}(M;\mathbb R)\)，它光滑依赖于 \(f\in H^s\)。记 \(Q=A^{-1}\)，证明在 \(H^{s+2}\) 范数中
	\[
		u=Qf-Q\bigl((Qf)^3\bigr)+O(\|f\|_{H^s}^5).
	\]
\end{exercise}
\begin{solution}
	\(A\) 的核为零，因为 \(Au=0\) 与 \(u\) 配对给出 \(\|u\|_2^2+\|\mathrm du\|_2^2=0\)。它形式自伴，故 Fredholm 选择给出满射，正交核后的椭圆估计使 \(Q:H^s\to H^{s+2}\) 有界。映射 \(u\mapsto Au+u^3\) 从 \(H^{s+2}\) 到 \(H^s\) 光滑，且零点微分为 \(A\)。Banach 逆函数定理给出所述局部唯一解及光滑依赖，特别是 \(\|u\|_{H^{s+2}}\le C\|f\|_{H^s}\)。

	令 \(v=Qf\)。方程写成 \(u=v-Qu^3\)，代数估计首先给出
	\(\|u-v\|_{H^{s+2}}\le C\|f\|_{H^s}^3\)。再用
	\(u^3-v^3=(u-v)(u^2+uv+v^2)\)，有
	\[
		\|u-v+Qv^3\|_{H^{s+2}}
		\le C\|u-v\|_{H^s}
		\bigl(\|u\|_{H^s}^2+\|v\|_{H^s}^2\bigr)
		\le C\|f\|_{H^s}^5.
	\]
	这里每次乘积都在 \(s>m/2\) 的 Sobolev 代数中估计，随后用 \(Q\) 获得两阶导数。
\end{solution}

\begin{exercise}\label{g3:ex:sobolev-sphere-chart}
	设 \(M^m\) 闭，\(k\ge1\)、\(1<p<\infty\)、\(kp>m\)，且 \(u_0\in W^{k,p}(M;\mathbb R^{n+1})\) 满足 \(|u_0|=1\)。无需假定 \(u_0\) 光滑，证明
	\[
		K_{u_0}=\{h\in W^{k,p}(M;\mathbb R^{n+1}):h\cdot u_0=0\}
	\]
	为有界投影的像，并用它为 \(u_0\) 附近的球面值映射给出显式光滑图。
\end{exercise}
\begin{solution}
	乘法估计使 \(P h=h-(h\cdot u_0)u_0\) 为有界算子，\(|u_0|=1\) 给出 \(P^2=P\)、\(\operatorname{im}P=K_{u_0}\)。所以 \(K_{u_0}\) 是闭 Banach 子空间。

	对 \(h\in K_{u_0}\)、\(\|h\|_\infty<1\)，定义
	\[
		\Psi(h)=h+\sqrt{1-|h|^2}\,u_0.
	\]
	由于 \(h\perp u_0\)，有 \(|\Psi(h)|^2=|h|^2+1-|h|^2=1\)，且 \(\Psi(h)\cdot u_0>0\)。复合定理和乘法估计证明 \(\Psi\) 光滑，其逆为 \(u\mapsto u-(u\cdot u_0)u_0\)，定义于球面值映射满足 \(u\cdot u_0>0\) 的开邻域上。严格正性在紧源上有统一下界，所以逆图确实得到 \(\|h\|_\infty<1\)。在零点，\(\mathrm D\Psi(0)h=h\)，故这一模型空间也就是该映射空间在 \(u_0\) 处的切空间。
\end{solution}

\begin{exercise}\label{g3:ex:sobolev-translation-loss}
	在 \(S^1=\mathbb R/(2\pi\mathbb Z)\) 上取整数 \(k\ge1\)。证明平移作用
	\(\mathcal R(t,u)(x)=u(x+t)\) 从 \(\mathbb R\times H^k(S^1)\) 到 \(H^k(S^1)\) 连续，但不是处处 \(C^1\)。具体证明：曲线 \(t\mapsto\mathcal R(t,u)\) 在零点按 \(H^k\) 范数可微，当且仅当 \(u\in H^{k+1}\)。
\end{exercise}
\begin{solution}
	写 \(u=\sum_{n\in\mathbb Z}u_ne^{\mathrm inx}\)。平移乘以 \(e^{\mathrm int}\)，故是 \(H^k\) 等距算子。由
	\(|e^{\mathrm int}-1|\le2\) 及加权平方可和，控制收敛给出每个固定 \(u\) 的强连续性。再用
	\[
		\|\mathcal R(t,u)-\mathcal R(t_0,u_0)\|_{H^k}
		\le\|u-u_0\|_{H^k}
		+\|\mathcal R(t-t_0,u_0)-u_0\|_{H^k}
	\]
	得到联合连续性。

	若 \(u\in H^{k+1}\)，差商的 Fourier 系数为 \((e^{\mathrm int}-1)u_n/t\)，逐项趋于 \(\mathrm in u_n\)，并由 \(|(e^{\mathrm int}-1)/t|\le|n|\) 控制，所以在 \(H^k\) 中趋于 \(u'\)。反之，若差商在 \(H^k\) 中收敛，则其每个 Fourier 系数必趋于 \(\mathrm in u_n\)，故极限为分布导数 \(u'\in H^k\)，即 \(u\in H^{k+1}\)。

	取实函数 \(u(x)=\sum_{n\ge1}n^{-k-1}\cos(nx)\)。其 \(H^k\) 范数平方与 \(\sum n^{-2}\) 可比，而 \(H^{k+1}\) 范数平方发散。因此平移轨道在这个点不可按 \(H^k\) 微分，联合映射不可能处处为 \(C^1\)。
\end{solution}

\begin{exercise}\label{g3:ex:torus-gauge-representative}
	在平坦环面 \(\mathbb T^2=\mathbb R^2/(2\pi\mathbb Z)^2\) 上令
	\(a=\cos x\,\mathrm dx+\sin x\,\mathrm dy\)。求唯一的零均值周期函数 \(\phi\)，使 \(a'=a+\mathrm d\phi\) 满足 \(\mathrm d^*a'=0\)。计算 \(\mathrm da'\)、\(\mathrm d^*\mathrm da'\) 及能量 \(\tfrac12\|\mathrm da'\|_2^2\)。再说明加上 \(c\,\mathrm dx+d\,\mathrm dy\) 后，为什么这些常数一形式不能总由周期函数的微分消去。
\end{exercise}
\begin{solution}
	平坦坐标中 \(\mathrm d^*(b_1\mathrm dx+b_2\mathrm dy)=-\partial_xb_1-\partial_yb_2\)，所以 \(\mathrm d^*a=\sin x\)。所需方程为 \(\Delta\phi=-\sin x\)，零均值解为 \(\phi=-\sin x\)，于是
	\[
		a'=\sin x\,\mathrm dy,\qquad
		\mathrm da'=\cos x\,\mathrm dx\wedge\mathrm dy,
		\mathrm d^*\mathrm da'=\sin x\,\mathrm dy.
	\]
	最后一个等式可直接分部积分：对二形式 \(f\,\mathrm dx\wedge\mathrm dy\)，\(\mathrm d^*\) 的坐标表达为 \((\partial_yf)\mathrm dx-(\partial_xf)\mathrm dy\)。能量为
	\[
		\frac12\int_0^{2\pi}\int_0^{2\pi}\cos^2x\,\mathrm dx\,\mathrm dy=\pi^2.
	\]
	若 \(\mathrm d\psi=c\,\mathrm dx+d\,\mathrm dy\) 且 \(\psi\) 周期，则沿两个坐标圆周积分分别得 \(2\pi c=0\)、\(2\pi d=0\)。故非零常数一形式不能由这种规范变化消去；它们是满足 \(\mathrm d=\mathrm d^*=0\) 的剩余调和方向。
\end{solution}

\begin{exercise}\label{g3:ex:inverse-radius-degeneration}
	对 \(\varepsilon>0\)，令 \(F_\varepsilon(x)=\varepsilon x+x^2\)。计算经过原点的局部逆，找出它的实值存在区间与原点导数。证明无法选择一个与 \(\varepsilon\downarrow0\) 无关的零点邻域，使该逆对其中所有右端存在，并将失败与定量逆函数定理中的常数联系起来。
\end{exercise}
\begin{solution}
	经过原点的分支是
	\[
		G_\varepsilon(y)=\frac{-\varepsilon+\sqrt{\varepsilon^2+4y}}2,
		\qquad y\ge-\frac{\varepsilon^2}4.
	\]
	它在 \(y>-\varepsilon^2/4\) 上光滑，\(G_\varepsilon(0)=0\)、\(G_\varepsilon'(0)=1/\varepsilon\)；在端点与另一根合并，逆的导数发散。函数 \(F_\varepsilon\) 的全局最小值就是 \(-\varepsilon^2/4\)，所以比这个值更小的右端根本没有实解。任意固定的零点开邻域都包含某个固定负数，它对足够小的 \(\varepsilon\) 位于该最小值以下，故不能统一取逆。

	在原点 \(A=\varepsilon\)，\(\|A^{-1}\|=1/\varepsilon\)，且
	\(1-A^{-1}F_\varepsilon'(x)=-2x/\varepsilon\)。若要求压缩常数至多 \(1/2\)，则可取的对称原像球半径须满足 \(r\le\varepsilon/4\)，相应定量保证的右端半径为 \(O(\varepsilon^2)\)。这与显式根的退化尺度一致。
\end{solution}
